\chapter[Spinful Fermions and Umklapp Phases]{Spinful Fermions, Spin-Charge Separation, and Umklapp / Gapped Phases}
\label{ch:spinful-umklapp}

When physical electrons possess spin-$\frac{1}{2}$ degrees of freedom, the single-particle
Hilbert space quadruples into four chiral fermion species: $(r, \sigma) \in \{R, L\} \times \{\uparrow, \downarrow\}$.
This chapter establishes the constructive mathematical theory of spinful $1+1$D systems,
resolving the exact nature of spin--charge separation, multi-target pairing sectors,
and the formation of interaction-induced energy gaps.

First, we prove that charge and spin currents decouple identically as operator distributions,
$[J_c(x), J_s(y)] = 0$, while $SU(2)$ spin-rotation symmetry pins the spin Luttinger parameter
to $K_s = 1$.  Second, we construct the four-species Klein 2-cocycle on $\mathbb Z^4$,
proving exact anticommutation across all four species.  Third, we analyze spin-singlet
pairing and show that the two spin components connect fundamentally different charge-sector targets.
Finally, we formalize Umklapp scattering at half-filling ($k_F = \pi/(2a)$): we prove by
exact phase cancellation that the Umklapp interaction acts purely in the charge sector as
$\cos(\sqrt{8\pi}\phi_c)$ while leaving the spin sector completely decoupled and gapless.
Under repulsive interactions ($K_c < 1$), the Umklapp operator becomes relevant with scaling
dimension $x_{\mathrm{umklapp}} = 2K_c < 2$, generating the Mott insulating charge gap.
The chapter concludes with a constructive Schrieffer--Wolff analysis of finite energy-shell
projections and effective low-energy Hamiltonians.

\section{Spinful Species, Raw Currents, and the Parity Sublattice}
\label{sec:ch21-spin-charge}

\subsection*{From two chiral branches to four fermion species}

For spinless fermions, the two chiral branches carry separate density
fluctuations.  Adding an electron spin label produces four branches, but it
also raises a question which is invisible if one starts immediately with
continuum ``charge'' and ``spin'' fields: \emph{which pairs of charge and spin
quantum numbers actually describe finite fermionic states?}  A second question
concerns the familiar claim of spin--charge separation.  The operator algebra
will separate into simple linear combinations, whereas commutation on a
finite energy budget and factorization of the state space require distinct
arguments.  We first formulate these questions entirely in the finite CAR
(carrier) representation, before using any oscillator interpretation.

Fix an integer $h\geq 1$, and write
\begin{equation}
  L=2h,\qquad
  \mathcal B_L=(-h,h]\cap\mathbb Z=\{1-h,\ldots,h\},\qquad
  \Lambda_L=\mathbb Z/L\mathbb Z.
  \label{eq:ch21-finite-band}
\end{equation}
Thus $L$ is the number of lattice sites and also the number of allowed
momenta per species.  We reserve $M\geq 1$ for a \emph{positive-mode cutoff}:
when a collection of modes is considered, $1\leq m\leq M$; an individual
signed density shift $m$ may instead be any integer.  The Fourier root and
orientation are fixed by
\begin{equation}
  \zeta=\exp(2\pi i/L),\qquad
  c_{\alpha,x}=\frac{1}{\sqrt L}
     \sum_{k\in\mathcal B_L}\zeta^{kx}c_{\alpha,k},
  \qquad x\in\Lambda_L.
  \label{eq:ch21-fourier-orientation}
\end{equation}
We use periodic boundary conditions.  The inverse transform carries
$\zeta^{-kx}$.  The finite \emph{non-wrapping} density modes below must not be
silently identified with the complete spatial Fourier coefficients of
$c_{\alpha,x}^{\dagger}c_{\alpha,x}$: the latter also contain momenta
wrapped across the ends of $\mathcal B_L$ (see Chapter~\ref{ch:density_modes}).

\begin{definition}[Ordered spinful CAR carrier]
\label{def:ch21-species}
Let
\begin{equation}
  \mathcal C=\{L,R\}\times\{\uparrow,\downarrow\}.
  \label{eq:ch21-species}
\end{equation}
The chiral label is $\eta\in\{L,R\}$ and the spin label is
$\sigma\in\{\uparrow,\downarrow\}$.  For explicit sign computations we
choose once and for all the total order
\begin{equation}
  (L,\uparrow)<(L,\downarrow)<(R,\uparrow)<(R,\downarrow),
  \label{eq:ch21-species-order}
\end{equation}
followed by increasing momentum within each species.  The one-particle mode
set is $\mathcal I=\mathcal C\times\mathcal B_L$.  The fermionic Hilbert
space has the orthonormal occupation basis
\begin{equation}
  \mathcal H_L=\ell^2\!\bigl(\mathcal P(\mathcal I)\bigr),
  \quad |S\rangle,\ S\subseteq\mathcal I,
  \qquad \dim\mathcal H_L=2^{4L}.
\end{equation}
For $i\in\mathcal I$ write $c_i$ and $c_i^\dagger$ for the
annihilation and creation maps.  With $\{A,B\}=AB+BA$ they obey
\begin{equation}
 \{c_i,c_j\}=\{c_i^\dagger,c_j^\dagger\}=0,
 \qquad \{c_i,c_j^\dagger\}=\delta_{ij}I.
 \label{eq:ch21-four-species-car}
\end{equation}
\end{definition}

The order in \eqref{eq:ch21-species-order} is not a spin anisotropy.  It is a
convention needed to assign signs to occupation kets, CAR words and, later,
Klein operators.  To see the mechanism, put
$\epsilon(i,S)=(-1)^{\#\{j\in S:j<i\}}$.  Then
\begin{align}
 c_i|S\rangle
   &=\begin{cases}\epsilon(i,S)|S\setminus\{i\}\rangle,&i\in S,\\
                     0,&i\notin S,\end{cases}\label{eq:ch21-ann-action}\\
 c_i^\dagger|S\rangle
   &=\begin{cases}\epsilon(i,S)|S\cup\{i\}\rangle,&i\notin S,\\
                     0,&i\in S.\end{cases}\label{eq:ch21-cre-action}
\end{align}
These formulas reproduce the CAR: for $i\ne j$ the two possible orders
of inserting or deleting distinct occupied modes have opposite signs,
whereas for $i=j$ exactly one of $c_ic_i^\dagger$ and
$c_i^\dagger c_i$ acts as the identity on each basis ket.
The construction and its general CAR proof were established in
Chapter~\ref{ch:car_fock}; its multispecies version and the canonical
occupation tensor coordinates are supplied by the Core A05 module.
This matters because fermion operators of different species
\emph{anticommute}, even though the corresponding even density operators
will commute.

Let $n_{\alpha,k}=c_{\alpha,k}^\dagger c_{\alpha,k}$ and define the
sea-subtracted species-charge observable
\begin{equation}
  \widehat N_\alpha=\sum_{k\in\mathcal B_L}n_{\alpha,k}-hI,
  \qquad \alpha\in\mathcal C.
  \label{eq:ch21-relative-species-charge}
\end{equation}
Here the reference sea occupies the $h$ momenta $k\leq 0$ in every
species.  Its relative charges vanish; an occupation configuration has
$N_\alpha\in[-h,h]\cap\mathbb Z$.  This subtraction is essential when
talking about zero modes: $\sum_k n_{\alpha,k}$ counts \emph{absolute}
particles, while $\widehat N_\alpha$ counts the integer charge relative
to the chosen sea.

For an admissible vector of species charges
$\mathbf N=(N_\alpha)_{\alpha\in\mathcal C}$ and an integer $K\geq0$, let
\begin{equation}
 \mathcal B_{\mathbf N,K}
  =\operatorname{span}\Bigl\{|S\rangle:
     N_\alpha(S)=N_\alpha\ \text{for all }\alpha,
     \ \mathcal E(S)=\sum_{\alpha\in\mathcal C}\mathcal E_\alpha(S)
          \leq K\Bigr\}.
 \label{eq:ch21-joint-budget}
\end{equation}
Each $\mathcal E_\alpha$ is the nonnegative, integer excitation grading
\emph{above the ground configuration of its own charge sector}, constructed
in Chapter~\ref{ch:vacuum_budget}.  Equation
\eqref{eq:ch21-joint-budget} defines a coordinate subspace, not an
operator-order inequality of the form $H\psi\leq K\psi$.
We write $P_{\mathbf N,K}$ for its orthogonal projection.  The budget
constrains the \emph{sum} of the four excitation energies; it does not
place a separate independent cutoff $K$ on every species.

\subsection*{Raw currents and their exact inverse}

For $m\in\mathbb Z$ and $\alpha\in\mathcal C$, define the ordered
non-wrapping density operator
\begin{equation}
  \rho_{\alpha,m}
   =\sum_{\substack{k\in\mathcal B_L\\k+m\in\mathcal B_L}}
        c_{\alpha,k+m}^{\dagger}c_{\alpha,k}.
  \label{eq:ch21-species-density}
\end{equation}
The sign convention is deliberate: $m>0$ raises the momentum and the
excitation grading, and
$\rho_{\alpha,m}^{\dagger}=\rho_{\alpha,-m}$.  Since
$\rho_{\alpha,m}$ is even in fermion operators, the global CAR imply
\begin{equation}
 [\rho_{\alpha,m},\rho_{\beta,n}]=0\quad(\alpha\ne\beta).
 \label{eq:ch21-cross-species-density}
\end{equation}
For example, moving a two-fermion word for $\rho_{\alpha,m}$ past
one for $\rho_{\beta,n}$ entails four minus signs.  No energy budget is
needed for \eqref{eq:ch21-cross-species-density}.

\begin{definition}[Unnormalized charge and spin currents]
\label{def:ch21-raw-currents}
At each chirality $\eta$ put
\begin{align}
 R^c_{\eta,m}&=\rho_{(\eta,\uparrow),m}
                 +\rho_{(\eta,\downarrow),m},
     \label{eq:ch21-raw-charge}\\
 R^s_{\eta,m}&=\rho_{(\eta,\uparrow),m}
                 -\rho_{(\eta,\downarrow),m}.
     \label{eq:ch21-raw-spin}
\end{align}
Here $R^{c,s}$ denotes a \emph{raw current}, not the right-moving
chirality $R$.  Neither current contains a factor $1/\sqrt2$.
\end{definition}

The change of variables is algebraically invertible over $\mathbb C$:
\begin{equation}
 \begin{pmatrix}R^c_{\eta,m}\\R^s_{\eta,m}\end{pmatrix}
 =\begin{pmatrix}1&1\\1&-1\end{pmatrix}
   \begin{pmatrix}\rho_{(\eta,\uparrow),m}\\
                   \rho_{(\eta,\downarrow),m}\end{pmatrix},
 \qquad
 \begin{pmatrix}\rho_{(\eta,\uparrow),m}\\
                   \rho_{(\eta,\downarrow),m}\end{pmatrix}
 =\frac12\begin{pmatrix}1&1\\1&-1\end{pmatrix}
   \begin{pmatrix}R^c_{\eta,m}\\R^s_{\eta,m}\end{pmatrix}.
 \label{eq:ch21-raw-inverse}
\end{equation}
Indeed, the square of the displayed integer matrix is $2I$.
Taking adjoints and using that its entries are real proves
$(R^{c,s}_{\eta,m})^\dagger=R^{c,s}_{\eta,-m}$.  The factors
$1/2$ in the inverse are harmless for operators over $\mathbb C$;
their meaning for \emph{integer charge eigenvalues} will be quite
different.

\subsection*{The edge obstruction and the finite common-margin theorem}

A tempting argument says that up- and down-spin densities commute and
therefore so do the charge and spin currents.  The conclusion does
\emph{not} follow.  The two same-spin commutators still remain, and a
finite band has edge-dependent corrections to the usual current algebra.
Define
\begin{equation}
 E_{\eta,\sigma}(m,n)
  =[\rho_{(\eta,\sigma),m},\rho_{(\eta,\sigma),n}].
 \label{eq:ch21-edge-operator}
\end{equation}
These are genuine finite-dimensional operators, not a priori scalar
multiples of the identity.

\begin{proposition}[Global mixed-current identity]
\label{prop:ch21-global-edge}
For every $m,n\in\mathbb Z$ and either chirality $\eta$,
\begin{align}
 [R^c_{\eta,m},R^s_{\eta,n}]
   &=E_{\eta,\uparrow}(m,n)-E_{\eta,\downarrow}(m,n),
       \label{eq:ch21-mixed-edge}\\
 [R^c_{\eta,m},R^c_{\eta,n}]
   &=[R^s_{\eta,m},R^s_{\eta,n}]
     =E_{\eta,\uparrow}(m,n)+E_{\eta,\downarrow}(m,n).
       \label{eq:ch21-self-edge}
\end{align}
If $\eta\ne\eta'$, every commutator between
$R^a_{\eta,m}$ and $R^b_{\eta',n}$ vanishes for
$a,b\in\{c,s\}$, as an identity on all of $\mathcal H_L$.
\end{proposition}
\begin{proof}
Use bilinearity of $[A,B]=AB-BA$.  For instance, suppressing
$\eta,m,n$ when unambiguous,
\begin{align*}
 [R^c_m,R^s_n]
  &=[\rho_{\uparrow,m}+\rho_{\downarrow,m},
      \rho_{\uparrow,n}-\rho_{\downarrow,n}]\\
  &=[\rho_{\uparrow,m},\rho_{\uparrow,n}]
    -[\rho_{\uparrow,m},\rho_{\downarrow,n}]
    +[\rho_{\downarrow,m},\rho_{\uparrow,n}]
    -[\rho_{\downarrow,m},\rho_{\downarrow,n}].
\end{align*}
The middle two terms vanish by
\eqref{eq:ch21-cross-species-density}, giving
\eqref{eq:ch21-mixed-edge}.  With two plus signs, or two minus signs,
the same expansion gives \eqref{eq:ch21-self-edge}.  For different
chiralities all four species in the expansion are distinct, so again
every constituent commutator vanishes.  No rearrangement of fermionic
\emph{odd} factors was used.
\end{proof}

The single-species Schwinger coefficient in the convention of this book is
\begin{equation}
 \kappa(m,n)=
 \begin{cases}-m,&m+n=0,\\0,&m+n\ne0.\end{cases}
 \label{eq:ch21-schwinger-coefficient}
\end{equation}
It is the coefficient of a commutator acting on an appropriate budget,
not a globally valid nonzero scalar CCR on the finite Fock space.  The
origin of this distinction is the exact diagonal edge formula and the
frozen-margin theorem of Chapters~\ref{ch:heisenberg_algebra} and
\ref{ch:vacuum_budget}; their multispecies current version is provided
in the Core A06 support module.

\begin{theorem}[Raw charge/spin current algebra on a common budget]
\label{thm:ch21-budget-ccr}
Assume $h>0$, $K\geq0$, and an admissible charge vector
$\mathbf N$.  Fix $\eta\in\{L,R\}$ and $m,n\in\mathbb Z$.
Suppose \emph{both} spin species at that chirality satisfy
\begin{equation}
 \begin{split}
 |m|+|n|+K+|N_{(\eta,\uparrow)}|&\leq h,\\
 |m|+|n|+K+|N_{(\eta,\downarrow)}|&\leq h.
 \end{split}
 \label{eq:ch21-common-margin}
\end{equation}
For every $v\in\mathcal B_{\mathbf N,K}$, with the commutators
formed in the \emph{ambient} finite CAR representation,
\begin{align}
 [R^c_{\eta,m},R^s_{\eta,n}]v&=0,
   \label{eq:ch21-cross-ccr-budget}\\
 [R^c_{\eta,m},R^c_{\eta,n}]v
   &=[R^s_{\eta,m},R^s_{\eta,n}]v
     =2\kappa(m,n)v.
   \label{eq:ch21-self-ccr-budget}
\end{align}
In particular, if $m\geq1$ and
$2m+K+\max_{\sigma}|N_{(\eta,\sigma)}|\leq h$, then
\begin{equation}
 [R^c_{\eta,-m},R^c_{\eta,m}]v
   =[R^s_{\eta,-m},R^s_{\eta,m}]v=2m\,v,
 \qquad [R^c_{\eta,-m},R^s_{\eta,m}]v=0.
 \label{eq:ch21-two-m}
\end{equation}
\end{theorem}
\begin{proof}
The established single-species signed-current theorem states that,
under the corresponding margin in \eqref{eq:ch21-common-margin},
\begin{equation*}
 E_{\eta,\sigma}(m,n)v=\kappa(m,n)v
 \qquad(\sigma=\uparrow,\downarrow).
\end{equation*}
The same vector $v$ belongs to the joint budget, so each of these
identities applies \emph{to that vector}; the two separate conditions
are needed because the relative species charges can differ.  Substitute
them into \eqref{eq:ch21-mixed-edge} to obtain
$(\kappa-\kappa)v=0$.  Substituting into
\eqref{eq:ch21-self-edge} gives $(\kappa+\kappa)v=2\kappa v$.
Finally, taking the ordered pair of signed shifts $(-m,m)$ gives
$\kappa(-m,m)=m$ and $|{-m}|+|m|=2m$.
This proves all statements without postulating a finite-dimensional
Heisenberg representation.
\end{proof}

For all $1\leq m,n\leq M$, the simple \emph{sufficient} uniform
condition
\begin{equation}
 2M+K+\max_{\substack{\eta\in\{L,R\}\\
                         \sigma\in\{\uparrow,\downarrow\}}}
                   |N_{(\eta,\sigma)}|\leq h
 \label{eq:ch21-uniform-margin}
\end{equation}
allows the theorem to be applied to every pair of shifts among
$\{-M,\ldots,M\}$ at both chiralities.  It is not claimed to be
necessary.  In particular, a proof involving a \emph{product of several
currents} must additionally account for the intermediate excursions;
checking a one-commutator bound does not certify a whole word.

\begin{example}[The first nontrivial charge/spin oscillator]
\label{ex:ch21-first-mode}
Take $h=2$, so $L=4$ and
$\mathcal B_L=\{-1,0,1,2\}$, with all species charges zero and
$K=0$.  The ground ket $\Omega$ occupies $-1$ and $0$ in each
species.  For a fixed chirality and spin, the only unblocked
$m=1$ hop from this ground configuration is $0\to1$; write
$|e_\sigma\rangle=\rho_{(\eta,\sigma),1}\Omega$.
The CAR sign formula gives
\begin{equation*}
 \rho_{(\eta,\sigma),-1}\Omega=0,
 \qquad \rho_{(\eta,\sigma),-1}|e_\sigma\rangle=\Omega,
 \qquad \langle e_\uparrow,e_\downarrow\rangle=0.
\end{equation*}
Thus $R^c_{\eta,1}\Omega=|e_\uparrow\rangle+|e_\downarrow\rangle$
whereas $R^s_{\eta,1}\Omega=|e_\uparrow\rangle-|e_\downarrow\rangle$.
Their squared norms are both $2$, and direct application gives
\begin{equation}
 [R^c_{\eta,-1},R^c_{\eta,1}]\Omega=2\Omega,
 \quad [R^s_{\eta,-1},R^s_{\eta,1}]\Omega=2\Omega,
 \quad [R^c_{\eta,-1},R^s_{\eta,1}]\Omega=0.
 \label{eq:ch21-example-two}
\end{equation}
The bounds hold because $2(1)+0+0\leq h$.  Notice that the
\emph{positive} shift creates an excitation in either spin species:
the relative plus or minus sign distinguishes charge and spin
combinations, while both retain the raw normalization $2m$.
\end{example}

The budget hypothesis is substantive.  For example, keep one chirality,
put its up-spin species into the \emph{completely filled} band, its
down-spin species into the sea ground, and leave the other two species
in their sea grounds.  Every nonzero-momentum density hop annihilates
the full up-spin band, whereas the down-spin sector obeys
$[\rho_{\downarrow,-1},\rho_{\downarrow,1}]\Omega_\downarrow
 =\Omega_\downarrow$ for $h\geq2$.  Consequently on the resulting
nonzero product ket $v$,
\begin{equation}
  [R^c_{\eta,-1},R^s_{\eta,1}]v=-v\ne0.
  \label{eq:ch21-off-margin-example}
\end{equation}
The up-spin relative charge is $h$, so the common bound fails.
This is a direct counterexample to unrestricted global spin/charge
commutation on the finite CAR carrier.

There is another important boundary effect even \emph{inside} the
regime of Theorem~\ref{thm:ch21-budget-ccr}.  Set
$P=P_{\mathbf N,K}$, $Q=I-P$, and let $A,B$ be any ambient maps.
Insertion of $I=P+Q$ gives the exact identity
\begin{equation}
 [PAP,PBP]
   =P[A,B]P-PAQBP+PBQAP.
 \label{eq:ch21-compression-residual}
\end{equation}
Here the left side contains \emph{compressed} operators, unlike the
ambient commutators of the theorem.  In the preceding $K=0$ example,
$P$ has rank one, $PR^{c}_{\eta,\pm1}P=0$ and therefore the
compressed commutator is zero, although
$P[R^c_{\eta,-1},R^c_{\eta,1}]P=2P$.  The difference is precisely
supplied by the excursions through $Q\mathcal H_L$.  This example
makes concrete why no finite-dimensional compression can realize a
nonzero scalar CCR as an identity on its entire range: taking the
trace of $[A,B]=\lambda I$ in finite dimension would give
$0=\lambda\dim\mathcal H$.

\subsection*{Integer zero modes and the parity sublattice}

The same linear transformation must now be applied to the integer
\emph{zero-mode} charges, not to oscillator operators.  For each
chirality, define the commuting self-adjoint observables
\begin{equation}
 \widehat Q_{c,\eta}=\widehat N_{(\eta,\uparrow)}
                         +\widehat N_{(\eta,\downarrow)},
 \qquad
 \widehat Q_{s,\eta}=\widehat N_{(\eta,\uparrow)}
                         -\widehat N_{(\eta,\downarrow)}.
 \label{eq:ch21-zero-modes}
\end{equation}
They are called charge and spin zero modes in \emph{raw integer units};
$Q_s$ is twice the $S^z$ polarization in the usual spin-$\tfrac12$
normalization.  On $\mathcal B_{\mathbf N,K}$ their eigenvalues are
\begin{equation}
 Q_{c,\eta}=N_{(\eta,\uparrow)}+N_{(\eta,\downarrow)},
 \qquad
 Q_{s,\eta}=N_{(\eta,\uparrow)}-N_{(\eta,\downarrow)}.
 \label{eq:ch21-zero-eigenvalues}
\end{equation}
This follows immediately from the diagonal species-charge observable
and its established action on a fixed sector.

\begin{proposition}[Parity sublattice and integral inverse]
\label{prop:ch21-parity-lattice}
Consider the map of integer charge lattices
\begin{equation}
 T:\mathbb Z^2\longrightarrow\mathbb Z^2,
 \qquad T(a,b)=(a+b,a-b).
 \label{eq:ch21-parity-map}
\end{equation}
Its image is exactly the index-two subgroup
\begin{equation}
 \Lambda_{c,s}
 =\{(q_c,q_s)\in\mathbb Z^2:q_c\equiv q_s\pmod2\}.
 \label{eq:ch21-parity-sublattice}
\end{equation}
On that subgroup $T$ is a bijection, with inverse
\begin{equation}
 T^{-1}(q_c,q_s)
 =\left(\frac{q_c+q_s}{2},\frac{q_c-q_s}{2}\right).
 \label{eq:ch21-parity-inverse}
\end{equation}
The inclusion $\Lambda_{c,s}\subsetneq\mathbb Z^2$ is proper;
for example, $(1,0)\notin\Lambda_{c,s}$.
\end{proposition}
\begin{proof}
For $a,b\in\mathbb Z$,
$T(a,b)_c-T(a,b)_s=(a+b)-(a-b)=2b$;
hence both coordinates have the same parity.  Conversely, assume
$q_c\equiv q_s\pmod2$.  Then $(q_c-q_s)/2$ is an integer, and
$q_c+q_s=(q_c-q_s)+2q_s$ is also even.  Therefore the two coordinates
in \eqref{eq:ch21-parity-inverse} are integers.  Applying $T$
to that pair yields
\begin{equation*}
  \left(\frac{q_c+q_s}{2}+\frac{q_c-q_s}{2},
        \frac{q_c+q_s}{2}-\frac{q_c-q_s}{2}\right)=(q_c,q_s).
\end{equation*}
Applying the proposed inverse to $T(a,b)$ returns $(a,b)$ in the
same way.  This proves injectivity, surjectivity onto the displayed
image, and the inverse formula.  Finally $1-0$ is odd, so $(1,0)$
is excluded.  Equivalently, the integer matrix for $T$ has determinant
$-2$; its image has index two and is not all of $\mathbb Z^2$.
\end{proof}

For the full four-species problem, the proposition applies
\emph{separately} at $\eta=L$ and $\eta=R$.  In other words, the
physical integer labels belong to
$\Lambda_{c,s}^{(L)}\times\Lambda_{c,s}^{(R)}$, with additional
finite-band charge bounds inherited from $-h\leq N_\alpha\leq h$.
A single parity statement for the \emph{sum} of the two chiralities
would be weaker and would not recover both independent species-charge
pairs.  The parity rule does not project out fermionic states: it is an
arithmetic description of the image of the existing species labels
under a non-unimodular change of coordinates.

\begin{example}[An allowed zero mode and a forbidden formal label]
Take $h=2$ and one chirality with
$N_\uparrow=2$ and $N_\downarrow=-1$.  The raw labels are
$(Q_c,Q_s)=(1,3)$.  They have equal parity, and
\eqref{eq:ch21-parity-inverse} recovers
$(N_\uparrow,N_\downarrow)=(2,-1)$.  By contrast, a hypothetical
$(Q_c,Q_s)=(1,0)$ would give
$(N_\uparrow,N_\downarrow)=(1/2,1/2)$; it is a perfectly reasonable
point of $\mathbb R^2$ but not a finite fermionic charge sector.
This is why replacing integer charges by two unconstrained
independent bosonic zero modes changes the state-label set.
\end{example}

\subsection*{What factorizes, and what does not}

Because the one-particle labels of different species are disjoint, the
occupation basis gives a unitary \emph{ordered} identification of the
unrestricted finite Hilbert space with the tensor product of the four
single-species occupation spaces.  For odd CAR operators the ordering
requires fermionic parity strings; for even density operators their
action can be taken on the corresponding tensor factor.  This
occupation-space factorization is a theorem of the finite CAR
construction, but it is \emph{not} a factorization into independent
charge and spin zero-mode lattices.

To make the energy restriction explicit, denote by
$\mathcal H_{\alpha,N_\alpha}^{(e)}$ the subspace of a single
species with relative charge $N_\alpha$ and excitation energy exactly
$e\in\mathbb N$.  Regrouping the occupation basis of
\eqref{eq:ch21-joint-budget} yields the canonical graded-vector-space
identification
\begin{equation}
 \mathcal B_{\mathbf N,K}\ \simeq\!
 \bigoplus_{\substack{(e_\alpha)_{\alpha\in\mathcal C}\in\mathbb N^4\\
                    \sum_\alpha e_\alpha\leq K}}
 \bigotimes_{\alpha\in\mathcal C}
       \mathcal H_{\alpha,N_\alpha}^{(e_\alpha)}.
 \label{eq:ch21-budget-energy-sum}
\end{equation}
The direct sum is orthogonal, and any zero-dimensional energy-grade
factor contributes nothing.  To prove the identity, map an occupation
ket to the ordered tensor product of its four restrictions to species.
This is a bijection of orthonormal bases; by definition of the total
excitation grading, the ket belongs to the budget precisely when the
four resulting grades have sum at most $K$.  Linearity proves the
claimed unitary identification.  No statement about charge/spin
\emph{oscillator} tensor factors is required for this proof.

In general,
\begin{equation}
  \mathcal B_{\mathbf N,K}
  \not\simeq_{\mathrm{natural}}
  \bigotimes_{\alpha\in\mathcal C}
      \Bigl(\bigoplus_{e=0}^{K}
             \mathcal H_{\alpha,N_\alpha}^{(e)}\Bigr),
 \label{eq:ch21-not-product-cutoff}
\end{equation}
where the notation emphasizes that the \emph{natural inclusion} of
the left-hand space into the right-hand tensor product is generally
strict; it is not a statement that abstract finite-dimensional vector
spaces could never have accidentally equal dimensions.  For example,
at $K=1$ a state with one energy-one excitation in each of two
species lies in the product of separate $K=1$ cuts, but has total
energy two and is absent from the joint budget.  In particular, a
putative $\mathcal B_{c,K}\otimes\mathcal B_{s,K}$ cannot be inferred
from a change of current variables, even before addressing the parity
sublattice.

\subsection*{Exchange-symmetric quadratic couplings}

We now isolate the precise sense in which spin and charge \emph{do}
separate algebraically.  Let $X_\uparrow,X_\downarrow$ and
$Y_\uparrow,Y_\downarrow$ be operators on a common vector space.
They may be density modes at different chiralities or momenta, and
we do \emph{not} assume that $X$ and $Y$ commute.  Introduce the raw
combinations $X_c=X_\uparrow+X_\downarrow$,
$X_s=X_\uparrow-X_\downarrow$, and similarly for $Y$.

\begin{proposition}[Ordered quadratic block identity]
\label{prop:ch21-quadratic-block}
For arbitrary $a,b\in\mathbb C$,
\begin{align}
 &a(X_\uparrow Y_\uparrow+X_\downarrow Y_\downarrow)
  +b(X_\uparrow Y_\downarrow+X_\downarrow Y_\uparrow)
 \notag\\
 &\hspace{1.5cm}=
  \frac{a+b}{2}\,X_cY_c+
  \frac{a-b}{2}\,X_sY_s.
 \label{eq:ch21-quadratic-block}
\end{align}
Thus a quadratic coupling matrix of the exchange-symmetric form
$\begin{psmallmatrix}a&b\\b&a\end{psmallmatrix}$ is block diagonal
in the raw charge/spin variables, with the coefficients displayed
above.  The statement preserves the product order throughout.
\end{proposition}
\begin{proof}
Expanding only by distributivity,
\begin{align*}
 X_cY_c&=X_\uparrow Y_\uparrow+X_\uparrow Y_\downarrow
          +X_\downarrow Y_\uparrow+X_\downarrow Y_\downarrow,\\
 X_sY_s&=X_\uparrow Y_\uparrow-X_\uparrow Y_\downarrow
          -X_\downarrow Y_\uparrow+X_\downarrow Y_\downarrow.
\end{align*}
The coefficients of the two same-spin terms become
$((a+b)+(a-b))/2=a$; those of the cross-spin terms become
$((a+b)-(a-b))/2=b$.  At no point do we commute an $X$ past a $Y$.
Equivalently, the change-of-basis matrix in
\eqref{eq:ch21-raw-inverse} diagonalizes
$\begin{psmallmatrix}a&b\\b&a\end{psmallmatrix}$ over
$\mathbb C$, but the explicit expansion shows that the matrix
argument remains valid for noncommutative operator entries.
\end{proof}

For a quadratic current Hamiltonian built as a sum of such ordered
bilinears, with an exchange-symmetric spin matrix \emph{for each}
choice of chirality and momentum labels, the proposition may be
applied term by term.  One obtains a charge-current part and a
spin-current part with no mixed bilinear.  Hermiticity of the
Hamiltonian requires the usual conjugate/reversed-index conditions
on its coefficients separately; spin exchange symmetry alone is not
an adjointness hypothesis.

The hypothesis is also necessary to inspect rather than assume.
For instance, a spin-dependent diagonal coupling with coefficients
$a_\uparrow\ne a_\downarrow$ becomes
\begin{equation}
 a_\uparrow X_\uparrow Y_\uparrow
 +a_\downarrow X_\downarrow Y_\downarrow
 =\frac{a_\uparrow+a_\downarrow}{4}(X_cY_c+X_sY_s)
  +\frac{a_\uparrow-a_\downarrow}{4}(X_cY_s+X_sY_c).
 \label{eq:ch21-asymmetry-mixes}
\end{equation}
The last two terms mix charge and spin.  Even when
\eqref{eq:ch21-quadratic-block} holds, it implies neither global
commutation of its finite-CAR charge and spin Hamiltonian parts nor a
unitary factorization of a \emph{compressed} Hamiltonian.  Global
current commutators retain the edge difference of
\eqref{eq:ch21-mixed-edge}; products require further excursion
bounds, and compressions retain the residual terms of
\eqref{eq:ch21-compression-residual}.  Nonquadratic interactions,
including the operators considered in later sections, are not
covered by this proposition.

\subsection*{Do the charge and spin modes have different velocities?}

A quadratic block decomposition permits distinct effective mode
velocities but does not \emph{force} them.  To state the condition
without presupposing a continuum limit, consider for each formal
charge/spin oscillator block $\nu\in\{c,s\}$ the stable two-chirality
pairing Hamiltonian
\begin{equation}
 H_{\nu,m}
 =m\Bigl[v_{1,\nu}
   (a_{R,\nu}^{\dagger}a_{R,\nu}
    +a_{L,\nu}^{\dagger}a_{L,\nu})
  +v_{2,\nu}(a_{R,\nu}^{\dagger}a_{L,\nu}^{\dagger}
                         +a_{R,\nu}a_{L,\nu})\Bigr],
 \quad m\geq1.
 \label{eq:ch21-pairing-hamiltonian}
\end{equation}
Here the $a$'s satisfy exact CCR on an \emph{abstract oscillator
carrier}; this is a model for the budget-stable current algebra, not
a global canonical representation by the finite-dimensional CAR
operators.  Assume real coefficients and
$v_{1,\nu}>|v_{2,\nu}|$.  Using the CCR gives
\begin{equation}
 \left[\begin{pmatrix}a_{R,\nu}\\a_{L,\nu}^{\dagger}\end{pmatrix},
                  H_{\nu,m}\right]
 =m\begin{pmatrix}
  v_{1,\nu}&v_{2,\nu}\\-v_{2,\nu}&-v_{1,\nu}
 \end{pmatrix}
 \begin{pmatrix}a_{R,\nu}\\a_{L,\nu}^{\dagger}\end{pmatrix}.
 \label{eq:ch21-bogoliubov-dynamic}
\end{equation}
The characteristic polynomial of this $2\times2$ matrix is
$\lambda^2-(v_{1,\nu}^2-v_{2,\nu}^2)$; therefore its positive
frequency, divided by $m$, is
\begin{equation}
 u_\nu=\sqrt{v_{1,\nu}^2-v_{2,\nu}^2}>0.
 \label{eq:ch21-block-speed}
\end{equation}
This is an exact calculation \emph{within the stated quadratic CCR
model}.  It follows that
\begin{equation}
 u_c\ne u_s
 \quad\Longleftrightarrow\quad
 v_{1,c}^2-v_{2,c}^2\ne v_{1,s}^2-v_{2,s}^2,
 \label{eq:ch21-speed-nondegeneracy}
\end{equation}
provided both radicands are strictly positive.  The equivalence
uses only injectivity of the positive square root.  Its reverse
implication would require re-examination without stability and
positivity.

If one parametrizes a block by real hyperbolic numbers
$c_\nu,s_\nu$ satisfying $c_\nu^2-s_\nu^2=1$, then one may use
\begin{equation}
 v_{1,\nu}=u_\nu(c_\nu^2+s_\nu^2),\qquad
 v_{2,\nu}=2u_\nu c_\nu s_\nu,\qquad
 g_\nu=(c_\nu-s_\nu)^2,\qquad
 g_\nu^{-1}=(c_\nu+s_\nu)^2.
 \label{eq:ch21-hyperbolic-parameters}
\end{equation}
The last identity follows by multiplication of the two squares and
$c_\nu^2-s_\nu^2=1$.  The symbol $g_\nu$ denotes a quadratic-model
Luttinger parameter, \emph{not} the real Umklapp coupling $g_U$
used later.  We also reserve $K$ exclusively for the finite energy
budget, not for the Luttinger parameter.

\begin{example}[Exchange symmetry without velocity separation]
\label{ex:ch21-equal-speeds}
Start with a spin-independent free linearized model of velocity
$v>0$.  Its spin-exchange-symmetric quadratic coupling has no
up/down mixing, $b=0$, and the same diagonal coefficient $a=v$
for both spins.  Equation~\eqref{eq:ch21-quadratic-block} gives
$\tfrac v2 X_cY_c+\tfrac v2 X_sY_s$; it \emph{does} have charge
and spin blocks.  After consistent canonical oscillator
normalization, both blocks have
\begin{equation*}
  v_{1,c}=v_{1,s}=v,\qquad v_{2,c}=v_{2,s}=0,
  \qquad u_c=u_s=v,
  \qquad c_c=c_s=1,\quad s_c=s_s=0.
\end{equation*}
Hence symmetry and even a quadratic charge/spin decomposition do not
imply unequal propagation speeds.  For contrast, the stable parameter
choices $(v_{1,c},v_{2,c})=(17,15)$ and
$(v_{1,s},v_{2,s})=(5,0)$ give $u_c=8$ and $u_s=5$.
The latter distinction comes from the \emph{parameters}, not merely
from the existence of the charge/spin basis.
\end{example}

\begin{remark}[Mathematical and formalization scope]
\label{rem:ch21-formalization-scope}
Three separate conclusions have been established in this section.
The four-species CAR and the global density identities concern the
\emph{entire finite Fock carrier}; the coefficient $2m$ and mixed
current cancellation concern the \emph{ambient action on specified
frozen-margin budget vectors}; and the dispersion computation
concerns a \emph{specified abstract quadratic CCR model}.  None
alone identifies its oscillator vacuum with a finite interacting
fermionic ground state, proves a thermodynamic limit, or provides a
spectral gap.  A charge/spin oscillator tensor-product description
would also have to respect the parity lattice and the total-energy
cutoff.  Such a construction is not inferred here.

The finite multispecies CAR and its support results are present in
the supplied Core modules, but the Chapter~21 spin/charge statements
as a combined Lean development are \emph{not yet implemented or
verified}.  The explicit conventions to reconcile at whole-book
integration are the momentum/Fourier orientation, sea subtraction for
zero modes, raw-versus-normalized current factors, ambient-versus-
compressed products, and the inherited sign order
\eqref{eq:ch21-species-order}.  For physical background on the
charge/spin description of one-dimensional fermions, see
\cite{giamarchi2004,gogolin1998,vondelft1998}; these references
are context rather than certificates for the finite margin theorem.
\end{remark}

\section{The Spin-Singlet Cooper Pair and Its Two Charge Channels}
\label{sec:ch21-singlet-pair}

\subsection*{Why one singlet operator has two different targets}

A spin-singlet Cooper pair is often represented by a single compact field expression in the continuum literature. However, on a finite fermionic Hilbert space, we must first ask a more elementary question: \emph{which occupation-number sectors does this operator connect?} When assembling a pair from a right mover and a left mover, the two antisymmetrized spin assignments remove particles from different species. Although they carry the same \emph{total} charge and spin, they do not possess the same four-component charge vector. Consequently, when their finite energy budgets are treated as typed linear maps, their operators have distinctly different target spaces. Furthermore, their relative fermionic signs must be explicitly computed and cannot simply be discarded under the heuristic phrase ``the Klein factors cancel.''

We continue with the ordered four-species CAR carrier of
Section~\ref{sec:ch21-spin-charge}.  Set
\begin{equation}
 L=2h,\quad h\in\mathbb N_{>0},\qquad
 \mathcal B_L=\{1-h,2-h,\ldots,h\},\qquad
 \Lambda_L=\mathbb Z/L\mathbb Z,\qquad
 \zeta=e^{2\pi i/L}.
 \label{eq:ch21-pair-setting}
\end{equation}
Every one of the four species
\(\mathcal C=\{L,R\}\times\{\uparrow,\downarrow\}\) has these \(L\)
momentum modes; the lexicographic occupation order is
\begin{equation}
 L\uparrow<L\downarrow<R\uparrow<R\downarrow,
 \quad\text{then increasing momentum within each species}.
 \label{eq:ch21-pair-order}
\end{equation}
We retain the positive-mode cutoff notation \(1\leq m\leq M\) for later
current or oscillator expansions, but the exact CAR identities in this
section require no such cutoff.  All site labels \(x\) are periodic lattice
positions, represented by integers modulo \(L\).  With periodic Fourier
orientation fixed as in Section~\ref{sec:ch21-spin-charge}, the finite site
annihilator is
\begin{equation}
 c_{\alpha,x}=\frac1{\sqrt L}\sum_{k\in\mathcal B_L}
        \zeta^{kx}c_{\alpha,k},\qquad
 c_{\alpha,x}^{\dagger}=\frac1{\sqrt L}\sum_{k\in\mathcal B_L}
        \zeta^{-kx}c_{\alpha,k}^{\dagger}.
 \label{eq:ch21-pair-site-field}
\end{equation}
These are genuine operators on the finite occupation Fock space
\(\mathcal H_L\), not formal vertex exponentials.  Orthogonality of the
finite characters and the momentum CAR give
\(\{c_{\alpha,x},c_{\beta,y}^{\dagger}\}
=\delta_{\alpha\beta}\delta_{xy}I\) and
\(\{c_{\alpha,x},c_{\beta,y}\}=0\).
Indeed, the first anticommutator is
\(L^{-1}\delta_{\alpha\beta}\sum_{k\in\mathcal B_L}
\zeta^{k(x-y)}I\), which is the stated periodic Kronecker delta.
There is no non-wrapping current approximation in this step.

\begin{definition}[Annihilating spin-singlet pair]
\label{def:ch21-singlet}
The ordered, spin-antisymmetric two-fermion operator is
\begin{equation}
 \begin{split}
 O_{SSC}(x)&=A_1(x)-A_2(x),\\
 A_1(x)&=c_{R\uparrow,x}\,c_{L\downarrow,x},\\
 A_2(x)&=c_{R\downarrow,x}\,c_{L\uparrow,x}.
 \end{split}
 \label{eq:ch21-singlet-def}
\end{equation}
The rightmost field acts first.  In particular, \(A_1\) annihilates an
\(L\downarrow\) particle before an \(R\uparrow\) particle; \(A_2\)
annihilates an \(L\uparrow\) particle before an \(R\downarrow\) particle.
\end{definition}

The minus sign in \eqref{eq:ch21-singlet-def} is the antisymmetric spin
contraction \(\sum_{\sigma,\tau}\epsilon_{\sigma\tau}
 c_{R\sigma,x}c_{L\tau,x}\), with
\(\epsilon_{\uparrow\downarrow}=1\) and
\(\epsilon_{\downarrow\uparrow}=-1\).  For \(U\in SU(2)\), the elementary
identity \(U^{T}\epsilon U=(\det U)\epsilon=\epsilon\) shows that this
contraction is a spin scalar when both spin labels transform by the same
matrix.  This representation-theoretic assertion does \emph{not} say that
the Hamiltonian possesses spin-rotation symmetry.  It also does not
eliminate the antisymmetric CAR ordering.  For example, taking the adjoint
and reordering the fermions gives
\begin{equation}
 O_{SSC}(x)^\dagger
  =-c_{R\uparrow,x}^\dagger c_{L\downarrow,x}^\dagger
    -c_{L\uparrow,x}^\dagger c_{R\downarrow,x}^\dagger.
 \label{eq:ch21-singlet-adjoint}
\end{equation}
Thus the annihilating convention here differs by an overall sign upon
adjunction from a frequently written spin-singlet \emph{creation} convention.

\subsection*{Four-component charge shifts and the common total quantum numbers}

Let \(\widehat N_\alpha=\sum_{k\in\mathcal B_L}
 c_{\alpha,k}^\dagger c_{\alpha,k}-hI\) be the relative species charge;
write \(\mathbf N\in\mathbb Z^{\mathcal C}\) for its simultaneous
eigenvalues, and \(e_\alpha\) for the unit charge vector in species
\(\alpha\).  Only admissible components \(-h\leq N_\alpha\leq h\)
occur on the finite carrier.  The \(c_{\alpha,k}\) preserve neither
\(\mathbf N\) nor the energy budget, but their \emph{charge shifts} are
exact on the entire ambient finite Fock space.

\begin{lemma}[Exact species shifts of the two singlet channels]
\label{lem:ch21-singlet-shifts}
For every \(\gamma\in\mathcal C\) and every \(x\in\Lambda_L\),
\begin{align}
 [\widehat N_\gamma,A_1(x)]
   &=-(\delta_{\gamma,R\uparrow}+\delta_{\gamma,L\downarrow})A_1(x),
   \label{eq:ch21-pair-shift-one}\\
 [\widehat N_\gamma,A_2(x)]
   &=-(\delta_{\gamma,R\downarrow}+\delta_{\gamma,L\uparrow})A_2(x).
   \label{eq:ch21-pair-shift-two}
\end{align}
Consequently, if \(\mathcal H_{\mathbf N}\subseteq\mathcal H_L\)
is the full joint eigenspace of species charges \(\mathbf N\), then
\begin{equation}
 A_j(x)\mathcal H_{\mathbf N}\subseteq\mathcal H_{\mathbf T_j},
 \qquad\begin{cases}
 \mathbf T_1=\mathbf N-e_{L\downarrow}-e_{R\uparrow},\\
 \mathbf T_2=\mathbf N-e_{L\uparrow}-e_{R\downarrow}.
 \end{cases}
 \label{eq:ch21-pair-targets}
\end{equation}
These two targets are distinct integer vectors.  Their charge eigenspaces
are orthogonal, including when one of the eigenspaces is zero.
\end{lemma}

\begin{proof}
The CAR imply
\([c_{\gamma,k}^\dagger c_{\gamma,k},c_{\beta,q}]
=-\delta_{\gamma\beta}\delta_{kq}c_{\beta,q}\).
Summing over \(k\), subtracting the scalar sea occupation \(hI\), and
then using \eqref{eq:ch21-pair-site-field} gives
\([\widehat N_\gamma,c_{\beta,x}]
=-\delta_{\gamma\beta}c_{\beta,x}\).
For any operators \(D,X,Y\), the identity
\([D,XY]=[D,X]Y+X[D,Y]\) holds without a commutativity
assumption.  Apply it to the \emph{ordered} words \(A_1,A_2\) to obtain
\eqref{eq:ch21-pair-shift-one}--\eqref{eq:ch21-pair-shift-two}.
If \(\widehat N_\gamma\psi=N_\gamma\psi\), then
\(\widehat N_\gamma A_j\psi=(N_\gamma+\Delta_{j,\gamma})A_j\psi\),
which gives \eqref{eq:ch21-pair-targets}.  Its two shift vectors are
different because, for instance, their \(L\uparrow\) components are
\(0\) and \(-1\).  Finally, distinct eigenspaces of a self-adjoint
operator are orthogonal: if \(\widehat N_\gamma u=au\) and
\(\widehat N_\gamma v=bv\), then
\(a\langle u,v\rangle=\langle\widehat N_\gamma u,v\rangle
=\langle u,\widehat N_\gamma v\rangle=b\langle u,v\rangle\).
Choose a component for which \(a\ne b\).
\end{proof}

It is instructive to express these shifts using the raw charge and spin
labels of Section~\ref{sec:ch21-spin-charge}:
\(Q_{\eta,c}=N_{\eta\uparrow}+N_{\eta\downarrow}\) and
\(Q_{\eta,s}=N_{\eta\uparrow}-N_{\eta\downarrow}\).
\begin{center}
\begin{tabular}{@{}lrrrrrr@{}}
\toprule
 & $\Delta Q_{R,c}$ & $\Delta Q_{R,s}$
 & $\Delta Q_{L,c}$ & $\Delta Q_{L,s}$
 & $\Delta Q_{\mathrm{tot},c}$ & $\Delta Q_{\mathrm{tot},s}$\\
\midrule
$A_1$ & $-1$ & $-1$ & $-1$ & $+1$ & $-2$ & $0$\\
$A_2$ & $-1$ & $+1$ & $-1$ & $-1$ & $-2$ & $0$\\
\bottomrule
\end{tabular}
\end{center}
The common total \(\Delta Q_{\mathrm{tot},c}=-2\) and
\(\Delta Q_{\mathrm{tot},s}=0\) do not identify the two targets.  In
particular, the chirality-resolved spin shifts differ, and each branch
continues to obey the integer parity constraint
\(Q_{\eta,c}\equiv Q_{\eta,s}\pmod 2\).  Forgetting chirality-resolved
charges before projecting would erase information needed to determine the
correct codomain.

\subsection*{The signed Klein factors are ordered charge-shift maps}

While a Klein factor is sometimes introduced as an abstract anticommuting symbol attached to each fermion field, on our finite occupation representation the source of its sign is entirely concrete: removing a fermion requires commuting it past all occupied modes that precede it in the fixed species--momentum order. Thus, the signed factor is a rigorously defined \emph{charge-dependent map}, rather than an arbitrary numerical phase assigned loosely to each species.

For each admissible integer vector \(\mathbf N\), write
\(\Omega_{\mathbf N}\) for the normalized charge-sector ground ket, with
species \(\alpha\) filled at
\(\{1-h,\ldots,N_\alpha\}\).  Its number of particles of species
\(\alpha\) is \(h+N_\alpha\).  Assume below that a named annihilation is
possible, namely \(N_\alpha\geq1-h\), and that all indicated lowered
charge vectors remain admissible.  In Chapter~\ref{ch:car_fock}'s
occupation action, the annihilator of the \emph{top occupied momentum}
\(k=N_\alpha\) has preceding-occupation exponent
\begin{equation}
 \begin{split}
 E_h(\mathbf N,\alpha)
   &=\sum_{\beta<\alpha}(h+N_\beta)
       +(h+N_\alpha-1),\\
 \kappa_h(\mathbf N,\alpha)
   &=(-1)^{E_h(\mathbf N,\alpha)}\in\{+1,-1\}.
 \end{split}
 \label{eq:ch21-pair-klein-phase}
\end{equation}
The first term counts occupied modes in preceding species; the second
counts those below the occupied top mode within species \(\alpha\).
The signed finite Klein kernel constructed in Core Chapter~13
reproduces precisely this phase on sector ground states, under the
charge and budget margins stated for that construction.

To keep the zero-mode notation transparent, denote by \(F_\alpha\) the
\emph{typed sector shift} whose ground-state action is
\begin{equation}
 F_\alpha\Omega_{\mathbf N}
  =\kappa_h(\mathbf N,\alpha)\,
    \Omega_{\mathbf N-e_\alpha}.
 \label{eq:ch21-pair-klein-action}
\end{equation}
This notation describes the Core Klein action on admissible ground
states; extending it to particular non-ground budgets requires the
corresponding finite kernel theorem.  It is not an assertion that one
bounded lowering operator acts isometrically on every charge sector,
including the empty boundary sector.  Define also the diagonal site phase
\(Z_{\alpha,x}\Omega_{\mathbf N}=\zeta^{xN_\alpha}
\Omega_{\mathbf N}\).  Since lowering species \(\beta\ne\alpha\)
does not change \(N_\alpha\), the ground-state actions of
\(Z_{\alpha,x}\) and \(F_\beta\) commute for distinct species.

\begin{proposition}[Two different signed zero-mode products]
\label{prop:ch21-two-klein-products}
With \(\mathbf T_1,\mathbf T_2\) as in
\eqref{eq:ch21-pair-targets}, the two ordered zero-mode combinations are
\begin{align}
 \mathcal K_1(x)&=
 F_{R\uparrow}F_{L\downarrow}
 Z_{R\uparrow,x}Z_{L\downarrow,x}:
  \mathbf N\longrightarrow\mathbf T_1,
 \label{eq:ch21-klein-one}\\
 \mathcal K_2(x)&=
 F_{R\downarrow}F_{L\uparrow}
 Z_{R\downarrow,x}Z_{L\uparrow,x}:
  \mathbf N\longrightarrow\mathbf T_2.
 \label{eq:ch21-klein-two}
\end{align}
On an admissible sector ground ket, the actions are
\begin{equation}
 \mathcal K_j(x)\Omega_{\mathbf N}
  =\zeta^{x(N_{a_j}+N_{b_j})}
    \kappa_h(\mathbf N-e_{b_j},a_j)
    \kappa_h(\mathbf N,b_j)\,\Omega_{\mathbf T_j},
 \label{eq:ch21-klein-channel-ground}
\end{equation}
where \((a_1,b_1)=(R\uparrow,L\downarrow)\) and
\((a_2,b_2)=(R\downarrow,L\uparrow)\).
The rightmost \(F_{b_j}\) acts first; in particular the second
Klein phase must be evaluated on \(\mathbf N-e_{b_j}\), not on
\(\mathbf N\).
\end{proposition}

\begin{proof}
Act with the rightmost diagonal \(Z\) operators on
\(\Omega_{\mathbf N}\); their product is
\(\zeta^{x(N_{a_j}+N_{b_j})}\).  Next \(F_{b_j}\) lowers
\(\mathbf N\) to \(\mathbf N-e_{b_j}\) with phase
\(\kappa_h(\mathbf N,b_j)\).  Finally \(F_{a_j}\) acts on that
\emph{intermediate} charge vector, giving
\(\kappa_h(\mathbf N-e_{b_j},a_j)\).
This proves \eqref{eq:ch21-klein-channel-ground}.
Notice that in our fixed ordering both \(b_1<a_1\) and
\(b_2<a_2\).  Hence
\(E_h(\mathbf N-e_b,a)=E_h(\mathbf N,a)-1\) whenever
\(b<a\) and \(a\ne b\).  Reversing the two shifts changes their
relative sign, which is the finite occupation-string origin of the
fermionic cross-species anticommutation.  No freedom to suppress that
sign has been introduced.
\end{proof}

The \(Z\) factors specify zero-mode phases, not the full oscillator
content of the site fields.  The formulas above are \emph{not} a
factorization theorem expressing \(c_{a,x}c_{b,x}\) on every budget
as a product of two compressed bosonic exponentials.  Such a statement
requires the single-field identifications, a common domain of composition,
and control of intermediate projection remainders.

\begin{lemma}[Matching the ordered CAR action on sector grounds]
\label{lem:ch21-pair-ground-action}
For distinct species \(a,b\), assume that the charge-sector ground
\(\Omega_{\mathbf N}\) and both lowerings are admissible.  Let
\(k_b=N_b\) and \(k_a=N_a\) denote the occupied top momenta in
\(\mathcal B_L\).  Then
\begin{equation}
 c_{a,k_a}c_{b,k_b}\Omega_{\mathbf N}
  =\kappa_h(\mathbf N-e_b,a)\kappa_h(\mathbf N,b)
      \Omega_{\mathbf N-e_b-e_a},
 \label{eq:ch21-pair-ordered-top}
\end{equation}
and, for the site fields,
\begin{equation}
 \big\langle\Omega_{\mathbf N-e_b-e_a},
  c_{a,x}c_{b,x}\Omega_{\mathbf N}\big\rangle
   =\frac{\zeta^{x(N_a+N_b)}}{L}
     \kappa_h(\mathbf N-e_b,a)\kappa_h(\mathbf N,b).
 \label{eq:ch21-pair-channel-coefficient}
\end{equation}
The right-hand side is \(L^{-1}\) times the corresponding
zero-mode matrix coefficient in \eqref{eq:ch21-klein-channel-ground}.
\end{lemma}

\begin{proof}
For a ground configuration, deleting its top occupied momentum
\(k_b=N_b\) leaves the ground configuration of charge
\(\mathbf N-e_b\); the occupation-basis annihilation rule gives the
phase \(\kappa_h(\mathbf N,b)\).  Deleting \(k_a=N_a\) from this
intermediate ground gives the second phase
\(\kappa_h(\mathbf N-e_b,a)\), proving
\eqref{eq:ch21-pair-ordered-top}.
To compute the matrix element for site annihilators, expand both
finite Fourier sums in \eqref{eq:ch21-pair-site-field}.
The target ket is the \emph{ground} configuration at
\(\mathbf N-e_a-e_b\).  Removing any momentum other than the
occupied top mode of species \(b\) leaves a different occupation
configuration, and removal of a non-top momentum of species \(a\)
likewise cannot produce the target ground.  Orthogonality of
occupation kets therefore selects the unique summand
\((k_a,k_b)=(N_a,N_b)\).  Its two Fourier amplitudes multiply to
\(L^{-1}\zeta^{x(N_a+N_b)}\), while its CAR sign is the product
already obtained.  No current-CCR or energy-margin hypothesis enters
this matrix coefficient; it is an exact finite-CAR statement.
\end{proof}

\subsection*{A finite sign calculation: both Klein phases matter}

\begin{example}[Four species, \(L=4\), and two opposite singlet coefficients]
\label{ex:ch21-pair-sign}
Take \(h=2\), \(L=4\), \(x=1\), \(\zeta=i\), and the zero-charge sea
\(\Omega_{\mathbf0}\).  Each of the four species occupies momenta
\(\{-1,0\}\), ordered as in \eqref{eq:ch21-pair-order}.
Both singlet channels remove the top momentum \(k=0\) from their
respective species.

For channel 1, annihilating the rightmost \(L\downarrow\) field
encounters \(2+1=3\) preceding occupied modes, hence contributes
\(-1\).  In the intermediate state, annihilating \(R\uparrow\) at
\(k=0\) encounters \(2+1+1=4\) preceding modes, hence contributes
\(+1\).  Thus
\begin{equation}
 \kappa_2(\mathbf0,L\downarrow)=-1,\quad
 \kappa_2(-e_{L\downarrow},R\uparrow)=+1,\quad
 \langle\Omega_{\mathbf T_1},A_1(1)\Omega_{\mathbf0}\rangle=-\tfrac14.
 \label{eq:ch21-pair-example-one}
\end{equation}
For channel 2, removing \(L\uparrow\) at \(k=0\) encounters one
preceding occupied mode and contributes \(-1\).  Removing
\(R\downarrow\) at \(k=0\) from the intermediate configuration
encounters \(1+2+2+1=6\) preceding modes and contributes \(+1\).
Therefore
\begin{equation}
 \kappa_2(\mathbf0,L\uparrow)=-1,\quad
 \kappa_2(-e_{L\uparrow},R\downarrow)=+1,\quad
 \langle\Omega_{\mathbf T_2},A_2(1)\Omega_{\mathbf0}\rangle=-\tfrac14.
 \label{eq:ch21-pair-example-two}
\end{equation}
The Fourier phase is \(\zeta^{(0+0)x}=1\) in both cases, and each
pair contributes the normalization \(1/L=1/4\).  In the \emph{singlet
operator} the second channel is subtracted.  Accordingly,
\begin{equation}
 \boxed{\quad
 \langle\Omega_{\mathbf T_1},O_{SSC}(1)\Omega_{\mathbf0}\rangle
   =-\tfrac14,\qquad
 \langle\Omega_{\mathbf T_2},O_{SSC}(1)\Omega_{\mathbf0}\rangle
   =+\tfrac14.\quad}
 \label{eq:ch21-pair-example-final}
\end{equation}
There is no contradiction in the opposite signs: the target kets
belong to different, mutually orthogonal charge sectors.  Their
amplitudes must not be summed as if they were two contributions to
the same basis coefficient.
\end{example}

The sign pattern has a useful exact generalization.  With
\(\mathbf N=\mathbf0\) and arbitrary \(h\geq1\), the channel-1
preceding exponents are \(2h-1\) for \(L\downarrow\) and
\(3h-2\) for \(R\uparrow\) after the first removal.  Channel 2 has
exponents \(h-1\) for \(L\uparrow\) and \(4h-2\) for
\(R\downarrow\) after the first removal.  Therefore
\begin{equation}
 \kappa_h(-e_{L\downarrow},R\uparrow)
 \kappa_h(\mathbf0,L\downarrow)
 =\kappa_h(-e_{L\uparrow},R\downarrow)
 \kappa_h(\mathbf0,L\uparrow)=(-1)^{h+1}.
 \label{eq:ch21-pair-general-sign}
\end{equation}
At \(\mathbf N=\mathbf0\) the ground-to-ground coefficients are
\(\pm(-1)^{h+1}/L\) in the two mutually orthogonal targets,
with the overall negative sign applied to channel 2.

\subsection*{Typed finite budgets, the direct-sum target, and exact residuals}

The finite-field product is an ambient endomorphism on \(\mathcal H_L\), meaning \(A_1-A_2\) is well defined there. Trouble arises only if each annihilator is first compressed to a map between \emph{different} charge-sector energy budgets and these maps are then subtracted as if their codomains agreed. The remedy is to rigorously specify each intermediate space, calculate its omitted component, and embed both final targets into the same encompassing Hilbert space.

Let \(\mathcal E(S)\) be the nonnegative integer grading relative to
the ground configuration of the charge sector of \(S\), as in
Chapter~\ref{ch:vacuum_budget}.  For admissible \(\mathbf N\) and
\(J\in\mathbb N\), define the coordinate budget
\begin{equation}
 V_{\mathbf N,J}
  =\operatorname{span}\{\,|S\rangle:\mathbf N(S)=\mathbf N,
                          \ \mathcal E(S)\leq J\,\}
  \subseteq\mathcal H_L.
 \label{eq:ch21-pair-budget}
\end{equation}
Let \(i_{\mathbf N,J}:V_{\mathbf N,J}\hookrightarrow\mathcal H_L\)
be inclusion, \(p_{\mathbf N,J}=i_{\mathbf N,J}^\dagger\) its
orthogonal component projection, and
\(P_{\mathbf N,J}=i_{\mathbf N,J}p_{\mathbf N,J}\) the ambient
orthogonal projector.  A cutoff \(J\) is an excitation-energy budget,
not a bound on the scalar product \(\langle\psi,H\psi\rangle\).

For \(j=1,2\), retain the ordered pair \((a_j,b_j)\) above and choose
\emph{independent} source, intermediate, and target cutoffs
\(J_s,J_{m,j},J_{t,j}\in\mathbb N\).  Define
\begin{equation}
 \begin{aligned}
 V_s&=V_{\mathbf N,J_s},\\
 \mathbf M_j&=\mathbf N-e_{b_j},
    &V_{m,j}&=V_{\mathbf M_j,J_{m,j}},\\
 \mathbf T_j&=\mathbf N-e_{b_j}-e_{a_j},
    &V_{t,j}&=V_{\mathbf T_j,J_{t,j}}.
 \end{aligned}
 \label{eq:ch21-pair-three-spaces}
\end{equation}
If any shifted charge vector is not admissible, its corresponding
budget is interpreted as the zero subspace of \(\mathcal H_L\).
Write \(i_s,p_{m,j},i_{m,j},p_{t,j}\) for the associated inclusion
and coordinate projections.  Finally form the \emph{orthogonal sum}
\begin{equation}
 W=V_{t,1}\oplus V_{t,2},\qquad
 j_j:V_{t,j}\longrightarrow W,\qquad
 p_W:\mathcal H_L\longrightarrow W.
 \label{eq:ch21-pair-common-target}
\end{equation}
Here \(W\) is identified isometrically with the subspace
\(V_{t,1}+V_{t,2}\subseteq\mathcal H_L\); the sum is orthogonal by
Lemma~\ref{lem:ch21-singlet-shifts}.  Thus each \(j_j\) is the
canonical inclusion of one target into the common codomain.

\begin{theorem}[Exact two-channel projected singlet identity]
\label{thm:ch21-pair-projected}
For every finite \(L=2h\), site \(x\), admissible source charge,
and chosen nonnegative cutoffs, define the typed field maps
\begin{equation}
 \begin{aligned}
 B_{b,j}&=p_{m,j}c_{b_j,x}i_s:
                  V_s\longrightarrow V_{m,j},\\
 B_{a,j}&=p_{t,j}c_{a_j,x}i_{m,j}:
                  V_{m,j}\longrightarrow V_{t,j},\\
 \mathcal R_j&=p_{t,j}c_{a_j,x}
       (I-P_{m,j})c_{b_j,x}i_s:
                  V_s\longrightarrow V_{t,j},\\
 P_{m,j}&=i_{m,j}p_{m,j}:
                  \mathcal H_L\longrightarrow\mathcal H_L.
 \end{aligned}
 \label{eq:ch21-pair-typed-maps}
\end{equation}
Then, as an equality of linear maps \(V_s\longrightarrow W\),
\begin{equation}
 \boxed{\displaystyle
 p_W O_{SSC}(x)i_s
  =j_1\!\left(B_{a,1}\circ B_{b,1}+\mathcal R_1\right)
   -j_2\!\left(B_{a,2}\circ B_{b,2}+\mathcal R_2\right).}
 \label{eq:ch21-pair-typed-identity}
\end{equation}
No intermediate or final cutoff inequality is required.  In
particular, \(\mathcal R_1\) and \(\mathcal R_2\) are separately
specified, and neither is silently set to zero.
\end{theorem}

\begin{proof}
Fix a channel \(j\).  The insertion of the ambient orthogonal
identity \(I=P_{m,j}+(I-P_{m,j})\) \emph{between the two site
annihilators} is exact:
\begin{equation}
 \begin{aligned}
 p_{t,j}c_{a_j,x}c_{b_j,x}i_s
 &=p_{t,j}c_{a_j,x}
       P_{m,j}c_{b_j,x}i_s
   +p_{t,j}c_{a_j,x}(I-P_{m,j})c_{b_j,x}i_s\\
 &=\bigl(p_{t,j}c_{a_j,x}i_{m,j}\bigr)
       \bigl(p_{m,j}c_{b_j,x}i_s\bigr)
      +\mathcal R_j\\
 &=B_{a,j}\circ B_{b,j}+\mathcal R_j.
 \end{aligned}
 \label{eq:ch21-pair-insert-projection}
\end{equation}
This equality holds inside \(\operatorname{Hom}(V_s,V_{t,j})\).
By Lemma~\ref{lem:ch21-singlet-shifts}, the image of
\(c_{a_j,x}c_{b_j,x}i_s\) lies entirely in the
\emph{full} charge sector \(\mathcal H_{\mathbf T_j}\).
Since \(V_{t,1}\) and \(V_{t,2}\) occupy different charge sectors,
projecting the channel-\(j\) image to \(W\) yields precisely
\(j_jp_{t,j}c_{a_j,x}c_{b_j,x}i_s\); the projection to the other
target is zero.  Inserting the two channel identities
\eqref{eq:ch21-pair-insert-projection} into
\(p_W(A_1-A_2)i_s\) proves
\eqref{eq:ch21-pair-typed-identity}.  The two terms are subtracted
\emph{after} composition with \(j_1\) and \(j_2\), so their
codomain is genuinely the same \(W\).  This completes the proof.
\end{proof}

The theorem is a finite-Fock identity, not an approximate bosonization
formula.  Both residuals represent concrete vectors that passed through
the \emph{complement} of the selected intermediate energy budget before
being annihilated by the second field.  In particular,
\begin{equation}
 \mathcal R_j=0
 \quad\Longleftrightarrow\quad
 p_{t,j}c_{a_j,x}(I-P_{m,j})c_{b_j,x}i_s=0.
 \label{eq:ch21-pair-residual-criterion}
\end{equation}
The stronger range condition
\((I-P_{m,j})c_{b_j,x}i_s=0\) suffices, but is not necessary:
the second field and final projection can remove all discarded
components even when the intermediate field leaves the budget.  A
single-field bosonization candidate from Core Chapter~14 may replace
\(B_{a,j}\) or \(B_{b,j}\) only after its \emph{individual}
finite-budget identity and margins have been established.  A verified product of independently
compressed candidates is not automatically the compressed physical
product: precisely the \(\mathcal R_j\) term measures the difference.

\begin{example}[A nonzero intermediate-projection residual]
\label{ex:ch21-pair-residual}
Return to \(h=2\), \(L=4\), \(x=1\), source charge
\(\mathbf N=\mathbf0\) and source budget \(J_s=0\).  In channel 1,
take \(J_{m,1}=0\) but \(J_{t,1}=1\).  Starting from the sea
\(\Omega_{\mathbf0}\), remove momentum \(k=-1\), rather than the
top momentum \(k=0\), in species \(L\downarrow\).  The resulting
species occupation is \(\{0\}\), whereas the ground occupation of
relative charge \(-1\) is \(\{-1\}\).  It therefore has one unit
of excitation energy and is annihilated by the intermediate projection
\(P_{m,1}\).

Next remove the top momentum \(k=0\) in species \(R\uparrow\).
The final occupation ket \(\Xi\) has charge
\(\mathbf T_1=-e_{L\downarrow}-e_{R\uparrow}\), one excitation
unit in \(L\downarrow\), and zero excitation in every other
species, so \(\Xi\in V_{t,1}\).  No other momentum pair
produces this exact configuration.  The preceding counts are \(2\)
for the first \(L\downarrow\) removal and \(4\) for the
\(R\uparrow\) removal; both CAR signs are positive.  The two
Fourier factors contribute
\(L^{-1}\zeta^{-1}=(-i)/4\).  Consequently,
\begin{equation}
 \left\langle\Xi,\mathcal R_1\Omega_{\mathbf0}\right\rangle
     =-\frac{i}{4}\ne0,\qquad
 \left\langle\Xi,(B_{a,1}\circ B_{b,1})
          \Omega_{\mathbf0}\right\rangle=0.
 \label{eq:ch21-pair-residual-witness}
\end{equation}
Thus even a source \emph{ground state} does not make a truncated
intermediate projection exact for an arbitrary final cutoff.  Taking
\(J_{m,1}=1\) retains the displayed intermediate ket, but that alone
is not a proof that the entire residual vanishes for all source
vectors and all final cutoffs.
\end{example}

\begin{remark}[An enlarged final codomain]
\label{rem:ch21-pair-enlarged-target}
Equation~\eqref{eq:ch21-pair-typed-identity} uses \(W\) \emph{equal to}
the indicated two-budget orthogonal sum.  If \(W'\) is any larger
closed subspace containing \(W\), with inclusion
\(\iota:W\hookrightarrow W'\), its orthogonal projection satisfies
\begin{equation}
 p_{W'}O_{SSC}(x)i_s
  =\iota\,p_WO_{SSC}(x)i_s
     +p_{W'}(I-P_W)O_{SSC}(x)i_s.
 \label{eq:ch21-pair-larger-target}
\end{equation}
The last term includes whatever part of either channel lies in the
larger target outside the two selected budgets; it must be retained.
It vanishes if \(W'=W\), or under an independently proved range
condition.  Thus merely saying that \(W'\) ``contains both target
budgets'' is insufficient to infer the same two-term formula.
\end{remark}

\subsection*{What is established, and what remains conditional}

In summary, we have derived three logically separate statements.  First,
\eqref{eq:ch21-pair-shift-one}--\eqref{eq:ch21-pair-targets} are exact
species-charge selection rules in the full finite CAR algebra.  They
show why the singlet has two, rather than one, typed charge channels.
Second, \eqref{eq:ch21-pair-klein-phase}--\eqref{eq:ch21-pair-general-sign}
compute their ordered fermionic phases on charge-sector ground kets;
these phases are the ground-state content of the finite signed Klein
construction.  Third, \eqref{eq:ch21-pair-typed-identity} is an exact
linear-map identity on arbitrary finite coordinate budgets, and
Example~\ref{ex:ch21-pair-residual} shows why its two residuals cannot
in general be omitted.

In a charge/spin collective-field description, one may substitute the
algebraic inverse currents
\(\rho_{\eta\uparrow}=(R^c_\eta+R^s_\eta)/2\) and
\(\rho_{\eta\downarrow}=(R^c_\eta-R^s_\eta)/2\) into candidate
oscillator phases.  That substitution is \emph{not} by itself a proof
that compressed exponentials factor into a single common charge factor
and a common spin factor.  The two singlet channels have different
Klein products, different charge-sector targets, and potentially
different intermediate-budget errors.  Any such rewriting must specify
the uncompressed algebra in which exponentials are defined, prove the
required commutation and ordering identities, and establish their
compatibility with projection.  The formal series of
Section~\ref{sec:ch20-formal-vertices} likewise cannot be evaluated
at a numerical parameter or identified with finite-CAR expectations
without further arguments.

No thermodynamic-limit correlation exponent, spectral gap, or
interacting-state identification follows from the operator identities
of this section.  The exact CAR results used here rest on the
previously constructed finite carrier and its supported Klein and
coordinate-projection results; the spinful Chapter~21 implementations
and the present manuscript remain to be independently formalized,
compiled against the full book, and reviewed.  In the next section,
Umklapp scattering changes four species charges at once, making the
same sector-typing and projection questions essential rather than
merely notational.

\section{Umklapp Scattering, Charge-Sector Leakage, and Finite Projection}
\label{sec:ch21-umklapp}

\subsection*{From a four-fermion process to a precise invariance question}

In the preceding sections the charge and spin currents were built from four
ordered fermion species, and the spin-singlet pair was found to have two
separate species-charge targets.  We now confront a related question for a
Hamiltonian rather than an order parameter: \emph{does a chosen finite
low-energy space remain invariant under a local interaction?}  An answer
requires more than the name ``Umklapp.''  It requires an ordered CAR word,
an explicit charge shift, and a state on which the operator acts nontrivially
\emph{outside} the chosen space.  Pauli exclusion can make a seemingly
outward process vanish; the opposite process may remain nonzero while staying
inside the permitted charge box.

The discussion is entirely on a finite lattice and uses the occupation
representation from Chapter~\ref{ch:car_fock}.  We prove exact selection rules,
self-adjointness, and a conditional failure of invariance.  These statements
do not imply a thermodynamic charge gap or establish continuum renormalization
group equations.  For physical motivation and conventional bosonized models,
see, for example, \cite{giamarchi2004,gogolin1998}; our claims below are
finite algebraic statements with their hypotheses displayed.

As in Sections~\ref{sec:ch21-spin-charge} and
\ref{sec:ch21-singlet-pair}, choose \(h\in\mathbb N_{>0}\) and set
\begin{equation}
  L=2h,\qquad \mathcal B_L=(-h,h]\cap\mathbb Z
       =\{1-h,\ldots,h\},\qquad
  \Lambda_L=\mathbb Z/L\mathbb Z,\qquad \zeta=e^{2\pi i/L}.
  \label{eq:ch21-u-band}
\end{equation}
The four species have the fixed occupation ordering
\begin{equation}
 L\uparrow<L\downarrow<R\uparrow<R\downarrow,
 \quad\text{with increasing momentum inside each species}.
 \label{eq:ch21-u-order}
\end{equation}
For \(\alpha\in\mathcal C=\{L,R\}\times\{\uparrow,\downarrow\}\),
let \(c_{\alpha,k}\) satisfy the finite canonical anticommutation relations
(CAR).  We use periodic position fields in the Fourier convention
\begin{equation}
 c_{\alpha,x}=\frac1{\sqrt L}\sum_{k\in\mathcal B_L}
                \zeta^{kx}c_{\alpha,k},\qquad
 c_{\alpha,x}^{\dagger}=\frac1{\sqrt L}\sum_{k\in\mathcal B_L}
                \zeta^{-kx}c_{\alpha,k}^{\dagger}.
 \label{eq:ch21-u-site-fields}
\end{equation}
The positive-mode cutoff, if density excitations are considered separately,
is denoted \(M\) (with \(1\leq m\leq M\)); no density-mode or
current-CCR cutoff enters the CAR calculation in this section.

\subsection*{The ordered Umklapp word and its Hermitian Hamiltonian}

\begin{definition}[Finite-site Umklapp interaction]
\label{def:ch21-u-word}
Define the ordered product
\begin{equation}
  O_U(x)=c_{L\uparrow,x}^{\dagger}
         c_{L\downarrow,x}^{\dagger}
         c_{R\downarrow,x}c_{R\uparrow,x},
       \qquad x\in\Lambda_L.
  \label{eq:ch21-u-word}
\end{equation}
For a \emph{real} coupling \(g_U\in\mathbb R\), set
\begin{equation}
 H_U=\frac{g_U}{L}\sum_{x\in\Lambda_L}
           \bigl(O_U(x)+O_U(x)^\dagger\bigr).
 \label{eq:ch21-u-hamiltonian}
\end{equation}
Products act from right to left: the right-up and right-down particles
are removed, in that order, before the left-down and left-up particles are
created.  The factor \(1/L\) in \eqref{eq:ch21-u-hamiltonian} is fixed
throughout, including the finite coefficient examples below.
\end{definition}

In the finite Fock representation these are bounded operators on the
same Hilbert space \(\mathcal H_L\), so no operator-domain problem arises.
The adjoint reverses \emph{the order} of the four factors:
\begin{equation}
 O_U(x)^\dagger
  =c_{R\uparrow,x}^\dagger c_{R\downarrow,x}^\dagger
    c_{L\downarrow,x}c_{L\uparrow,x}.
 \label{eq:ch21-u-adjoint}
\end{equation}
It is not permissible to reorder these odd operators without keeping
the resulting CAR signs.

\begin{proposition}[Self-adjointness]
\label{prop:ch21-u-selfadjoint}
For every \(h>0\) and real \(g_U\), the operator \(H_U\) in
\eqref{eq:ch21-u-hamiltonian} satisfies \(H_U^\dagger=H_U\).
\end{proposition}
\begin{proof}
Write \(A=\sum_xO_U(x)\).  Since the sum is finite,
\(A^\dagger=\sum_xO_U(x)^\dagger\); hence
\[
 H_U^\dagger
 =\overline{\frac{g_U}{L}}(A^\dagger+A)
 =\frac{g_U}{L}(A+A^\dagger)=H_U.
\]
The middle equality uses the reality of \(g_U\).  No commutation
between \(A\) and \(A^\dagger\) was assumed.
\end{proof}

An Umklapp word is a possible lattice process rather than a statement of
its nonzero action on all states.  Its momentum-space expansion follows
exactly from \eqref{eq:ch21-u-site-fields}:
\begin{align}
 O_U(x)=\frac1{L^2}\!
 \sum_{k_1,k_2,k_3,k_4\in\mathcal B_L}
 &\zeta^{(-k_1-k_2+k_3+k_4)x}\,
 c_{L\uparrow,k_1}^\dagger c_{L\downarrow,k_2}^\dagger
 c_{R\downarrow,k_3}c_{R\uparrow,k_4}.
 \label{eq:ch21-u-fourier-expansion}
\end{align}
Summation over \(x\) retains precisely those ordered words satisfying
\(k_1+k_2-k_3-k_4\equiv0\pmod L\).  Thus momentum is conserved
\emph{modulo} the reciprocal-lattice period, rather than as an
unrestricted equality between the chosen integer representatives.
This distinction will be essential for the explicit witness.

\subsection*{Charge shifts and the surviving conservation laws}

The relative species number operators are
\begin{equation}
 \widehat N_\alpha=\sum_{k\in\mathcal B_L}
     c_{\alpha,k}^\dagger c_{\alpha,k}-hI,
 \qquad \alpha\in\mathcal C.
 \label{eq:ch21-u-relative-charge}
\end{equation}
They have integer eigenvalues \(-h,\ldots,h\) on the finite Fock space.
Let \(e_\alpha\) denote a basis vector in the four-dimensional
species-charge lattice.  With the order \eqref{eq:ch21-u-order} we define
\begin{equation}
 \Delta\mathbf N
  =e_{L\uparrow}+e_{L\downarrow}
     -e_{R\uparrow}-e_{R\downarrow}
  =(+1,+1,-1,-1).
 \label{eq:ch21-u-charge-vector}
\end{equation}
Every \emph{species} shift has magnitude one, whereas each chiral
\emph{branch total} changes by two.

\begin{lemma}[One-particle charge commutators]
\label{lem:ch21-u-fundamental-commutators}
For any species \(\gamma,\beta\), momentum \(q\in\mathcal B_L\), and
site \(x\),
\begin{align}
 [\widehat N_\gamma,c_{\beta,q}]
   &=-\delta_{\gamma\beta}c_{\beta,q},&
 [\widehat N_\gamma,c_{\beta,q}^\dagger]
   &=\delta_{\gamma\beta}c_{\beta,q}^\dagger,
 \label{eq:ch21-u-elementary-commutators}\\
 [\widehat N_\gamma,c_{\beta,x}]
   &=-\delta_{\gamma\beta}c_{\beta,x},&
 [\widehat N_\gamma,c_{\beta,x}^\dagger]
   &=\delta_{\gamma\beta}c_{\beta,x}^\dagger.
 \label{eq:ch21-u-site-commutators}
\end{align}
\end{lemma}
\begin{proof}
The CAR imply, for individual modes \(i,j\),
\[
 [c_i^\dagger c_i,c_j]
   =c_i^\dagger\{c_i,c_j\}-\{c_i^\dagger,c_j\}c_i
   =-\delta_{ij}c_j,
\]
and similarly \([c_i^\dagger c_i,c_j^\dagger]=\delta_{ij}c_j^\dagger\).
Summing over the momenta of species \(\gamma\) proves
\eqref{eq:ch21-u-elementary-commutators}; the scalar sea subtraction
commutes with everything.  Fourier linearity gives
\eqref{eq:ch21-u-site-commutators}.
\end{proof}

\begin{theorem}[Exact four-species Umklapp selection rule]
\label{thm:ch21-u-charge-shift}
For every \(x\) and species \(\gamma\),
\begin{align}
 [\widehat N_\gamma,O_U(x)]
   &=\Delta N_\gamma\,O_U(x),
 \label{eq:ch21-u-commutator}\\
 [\widehat N_\gamma,O_U(x)^\dagger]
   &=-\Delta N_\gamma\,O_U(x)^\dagger.
 \label{eq:ch21-u-adjoint-commutator}
\end{align}
Consequently, if \(\mathcal H_{\mathbf N}\) is the simultaneous
relative-charge eigenspace of \(\mathbf N\in\mathbb Z^{\mathcal C}\), then
\begin{equation}
 O_U(x)\mathcal H_{\mathbf N}\subseteq
     \mathcal H_{\mathbf N+\Delta\mathbf N},\qquad
 O_U(x)^\dagger\mathcal H_{\mathbf N}\subseteq
     \mathcal H_{\mathbf N-\Delta\mathbf N}.
 \label{eq:ch21-u-charge-targets}
\end{equation}
If a listed charge is outside the allowed finite range, its eigenspace
is understood to be zero.  The total relative particle number and the
spin imbalance satisfy
\begin{equation}
 \left[\sum_{\alpha\in\mathcal C}\widehat N_\alpha,H_U\right]=0,
 \qquad
 [\widehat S,H_U]=0,
 \quad \widehat S=\sum_{\eta=L,R}
        (\widehat N_{\eta\uparrow}-\widehat N_{\eta\downarrow}).
 \label{eq:ch21-u-conserved-totals}
\end{equation}
By contrast, \(H_U\) need not conserve each species charge or the
separate left/right branch charge.
\end{theorem}
\begin{proof}
For arbitrary operators, the commutator is a derivation:
\([D,AB]=[D,A]B+A[D,B]\).  Expanding the ordered word gives
\begin{align*}
 [\widehat N_\gamma,O_U(x)]={}&
   [\widehat N_\gamma,c_{L\uparrow,x}^\dagger]
           c_{L\downarrow,x}^\dagger c_{R\downarrow,x}c_{R\uparrow,x}
 +c_{L\uparrow,x}^\dagger
   [\widehat N_\gamma,c_{L\downarrow,x}^\dagger]
           c_{R\downarrow,x}c_{R\uparrow,x}\\
 &+c_{L\uparrow,x}^\dagger c_{L\downarrow,x}^\dagger
   [\widehat N_\gamma,c_{R\downarrow,x}]c_{R\uparrow,x}
 +c_{L\uparrow,x}^\dagger c_{L\downarrow,x}^\dagger
   c_{R\downarrow,x}[\widehat N_\gamma,c_{R\uparrow,x}].
\end{align*}
Lemma~\ref{lem:ch21-u-fundamental-commutators} reduces the coefficient
to
\(\delta_{\gamma,L\uparrow}+\delta_{\gamma,L\downarrow}
-\delta_{\gamma,R\downarrow}-\delta_{\gamma,R\uparrow}\), which is
exactly \(\Delta N_\gamma\).  Taking adjoints of the operator
commutator, using self-adjointness of \(\widehat N_\gamma\), gives
\eqref{eq:ch21-u-adjoint-commutator}.  For
\(\widehat N_\gamma\psi=N_\gamma\psi\),
\[
 \widehat N_\gamma O_U(x)\psi
  =(N_\gamma+\Delta N_\gamma)O_U(x)\psi,
\]
proving the first inclusion in \eqref{eq:ch21-u-charge-targets};
the second is identical with the opposite shift.  Finally the
coefficients of the total-number commutator sum to \(1+1-1-1=0\),
while those of \(\widehat S\) sum to \((1-1)+(-1+1)=0\).
The identities persist under the finite site sum and the adjoint.
\end{proof}

Writing \(\widehat Q_c^\eta=\widehat N_{\eta\uparrow}
 +\widehat N_{\eta\downarrow}\) and
\(\widehat Q_s^\eta=\widehat N_{\eta\uparrow}
 -\widehat N_{\eta\downarrow}\) as in \S\,21.1, the same computation
immediately gives
\begin{equation}
 [\widehat Q_c^L,O_U]=2O_U,\qquad
 [\widehat Q_c^R,O_U]=-2O_U,\qquad
 [\widehat Q_s^L,O_U]=[\widehat Q_s^R,O_U]=0.
 \label{eq:ch21-u-branch-shifts}
\end{equation}
Here \(O_U\) denotes either the fixed-site word or its outward site sum.
These equations describe an interbranch transfer of two particles,
not the creation or destruction of total electric charge.

\begin{corollary}[Vanishing same-sector compression]
\label{cor:ch21-u-sector-compression}
For any admissible charge vector \(\mathbf N\) and any two nonnegative
energy cutoffs \(K,K'\),
\begin{equation}
 P_{\mathbf N,K}\,H_U\,P_{\mathbf N,K'}=0.
 \label{eq:ch21-u-zero-compression}
\end{equation}
\end{corollary}
\begin{proof}
Both projected source and projected target are contained in
\(\mathcal H_{\mathbf N}\).  Theorem~\ref{thm:ch21-u-charge-shift}
places the two summands of \(H_U\) in
\(\mathcal H_{\mathbf N+\Delta\mathbf N}\) and
\(\mathcal H_{\mathbf N-\Delta\mathbf N}\), respectively.  Neither
charge vector equals \(\mathbf N\), and joint eigenspaces with
different charge labels are orthogonal.  The left projection therefore
kills both contributions, independently of \(K,K'\).
\end{proof}

\subsection*{Charge boxes, projections, and a conditional leakage theorem}

A single-charge-sector projection is not the same as a charge
\emph{box}, which admits several sectors but limits their extent.
To define the latter without using an informal inequality between
vectors, let \(S\subseteq\mathcal C\times\mathcal B_L\) be an
occupation configuration.  Its relative charge vector is
\(\mathbf N(S)\), and its nonnegative integer excitation energy
\(\mathcal E(S)\) is measured above the ground configuration of the
\emph{same} charge vector, as in Chapter~\ref{ch:vacuum_budget}.

\begin{definition}[Per-species charge box and budget]
\label{def:ch21-u-charge-box}
Given bounds \(B_\alpha\in\mathbb N\) and an energy budget
\(K\in\mathbb N\), let
\begin{equation}
 \mathcal B_{K,\mathbf B}
   =\operatorname{span}\Bigl\{|S\rangle:\;
       |N_\alpha(S)|\le B_\alpha\ \text{for all }\alpha,
       \quad\mathcal E(S)\le K\Bigr\}
       \subseteq\mathcal H_L,
 \label{eq:ch21-u-box}
\end{equation}
and let \(P_{K,\mathbf B}\) be its orthogonal coordinate projection.
The charge bounds and the energy cutoff play logically distinct roles.
A box may be proper even when the energy cutoff is large enough to
include every state whose charge vector satisfies the bounds.
\end{definition}

\begin{theorem}[Conditional outward leakage]
\label{thm:ch21-u-leakage}
Assume \(h\ge1\), \(g_U\in\mathbb R\setminus\{0\}\), and
\(\psi\in\mathcal B_{K,\mathbf B}\) has definite charge
\(\mathbf N\), i.e.\ \(\widehat N_\alpha\psi=N_\alpha\psi\)
for all \(\alpha\).  Suppose
\begin{equation}
  \sum_{x\in\Lambda_L}O_U(x)\psi\ne0,
  \qquad
  \exists\alpha\in\mathcal C:\quad
      |N_\alpha+\Delta N_\alpha|>B_\alpha.
 \label{eq:ch21-u-leakage-hypotheses}
\end{equation}
Then \(H_U\psi\notin\mathcal B_{K,\mathbf B}\), and the ambient
Hamiltonian does not preserve this budget.  In addition,
\begin{equation}
 [P_{K,\mathbf B},H_U]\ne0.
 \label{eq:ch21-u-projector-noncommute}
\end{equation}
An analogous conclusion holds for a nonzero adjoint (inward) site sum
whose target \(\mathbf N-\Delta\mathbf N\) violates a charge bound.
\end{theorem}
\begin{proof}
For any \(\mathbf T\), denote the orthogonal full-charge-sector
projection by \(\Pi_{\mathbf T}\).  Distinct sectors are orthogonal
because the \(\widehat N_\alpha\) are commuting self-adjoint matrices.
The two shifts in \eqref{eq:ch21-u-charge-targets} differ, hence
\begin{equation}
 \Pi_{\mathbf N+\Delta\mathbf N}H_U\psi
   =\frac{g_U}{L}\sum_xO_U(x)\psi\ne0.
 \label{eq:ch21-u-projected-outward}
\end{equation}
The adjoint contribution has charge \(\mathbf N-\Delta\mathbf N\)
and cannot cancel the displayed component.  By the second hypothesis,
\(\mathcal B_{K,\mathbf B}\) has no vector in the charge sector
\(\mathbf N+\Delta\mathbf N\).  Equation
\eqref{eq:ch21-u-projected-outward} therefore exhibits a nonzero
component of \((I-P_{K,\mathbf B})H_U\psi\).
Since \(P_{K,\mathbf B}\psi=\psi\), we also have
\[
 [P_{K,\mathbf B},H_U]\psi
  =P_{K,\mathbf B}H_U\psi-H_UP_{K,\mathbf B}\psi
  =-(I-P_{K,\mathbf B})H_U\psi\ne0.
\]
The proof for the opposite shift exchanges \(O_U\) with \(O_U^\dagger\).
\end{proof}

The same elementary argument detects \emph{energy} leakage: it suffices
to exhibit an occupation ket \(|T\rangle\) excluded by the energy
cutoff for which \(\langle T,H_U\psi\rangle\ne0\).
The charge-box hypothesis in Theorem~\ref{thm:ch21-u-leakage} is a
convenient sufficient condition, not a necessary condition for
non-invariance.  Conversely, \(H_U\psi\ne0\) by itself is
\emph{not} a leakage certificate.  The latter distinction will be
visible in the second finite example.

\subsection*{An explicit outward witness: the sign and the Fourier coefficient}

The next calculation checks every part of the leakage hypothesis with
a concrete occupation configuration.  It also illustrates why a
fermionic sign cannot be inferred from charge conservation alone.

\begin{theorem}[Zero-charge sea leaves every zero-charge box]
\label{thm:ch21-u-zero-box-witness}
For every \(h\ge1\), write \(L=2h\).  In each species fill the sea
momenta \(\{1-h,\ldots,0\}\), and let \(\Omega\) be the corresponding
normalized, ordered occupation ket.  Define \(T\) by removing
momentum \(1-h\) from \(R\uparrow\) and \(R\downarrow\), and
inserting momentum \(1\) in \(L\uparrow\) and \(L\downarrow\).
Then
\begin{equation}
 \langle T,H_U\Omega\rangle=-\frac{g_U}{L^2}.
 \label{eq:ch21-u-witness-coefficient}
\end{equation}
If \(g_U\ne0\), the zero-charge sea \(\Omega\) belongs to
\(\mathcal B_{0,\mathbf0}\), but
\(H_U\Omega\notin\mathcal B_{K,\mathbf0}\) for \emph{every}
\(K\ge0\).  Thus \(H_U\) does not preserve the zero-charge budget
and does not commute with its projector.
\end{theorem}
\begin{proof}
Let \(S_0\) be the set of \(4h\) occupied species-momentum modes
constituting \(\Omega\).  The target configuration is
\[
 T=\bigl(S_0\setminus\{(R\uparrow,1-h),(R\downarrow,1-h)\}\bigr)
     \cup\{(L\uparrow,1),(L\downarrow,1)\}.
\]
For \(h\ge1\), the removed momenta lie in the sea and the inserted
momentum \(1\) lies in the band but outside the sea.  Therefore no
Pauli factor kills this particular word.  The relative-charge vector
of \(T\) is \(\Delta\mathbf N=(1,1,-1,-1)\), whereas
\(\mathbf N(\Omega)=\mathbf0\).

For an occupied mode \(i\), annihilation contributes the sign
\((-1)^{\#\{j\in S:j<i\}}\), and creation at an unoccupied mode
contributes the corresponding sign of the updated configuration.
The ordered word \eqref{eq:ch21-u-word} acts from right to left.
The number of occupied modes \emph{preceding} each operation, after
all previous operations have been made, is
\begin{center}
\begin{tabular}{@{}lll@{}}
\toprule
Operation & Preceding occupied modes & Sign \\
\midrule
\(c_{R\uparrow,1-h}\) & \(2h\) & \((-1)^{2h}\) \\
\(c_{R\downarrow,1-h}\) & \(3h-1\) & \((-1)^{3h-1}\) \\
\(c_{L\downarrow,1}^\dagger\) & \(2h\) & \((-1)^{2h}\) \\
\(c_{L\uparrow,1}^\dagger\) & \(h\) & \((-1)^h\) \\
\bottomrule
\end{tabular}
\end{center}
For example, the first removal has the two filled left species before
it; after that removal the right-down species has \(2h+(h-1)\)
occupied modes before its lowest sea momentum.  The insertion in
left-down sees all \(h\) left-up and its own \(h\) filled sea modes,
while the final insertion in left-up sees its original \(h\) sea modes.
The exponent is \(2h+(3h-1)+2h+h=8h-1\), an odd integer.
Thus the selected momentum word produces \(-|T\rangle\).

In \eqref{eq:ch21-u-fourier-expansion}, the four momenta are
\(k_1=k_2=1\) and \(k_3=k_4=1-h\).  Their Fourier phase is
\[
 \zeta^{(-1-1+(1-h)+(1-h))x}
    =\zeta^{-2hx}=\zeta^{-Lx}=1
       \quad(x\in\Lambda_L).
\]
Each site word has scalar Fourier normalization \(L^{-2}\).
The four changed occupations uniquely fix the four participating
momenta: another term in the momentum sum cannot produce \(|T\rangle\)
from \(|\Omega\rangle\).  The adjoint word changes the charge by
\(-\Delta\mathbf N\), hence also cannot contribute to this target.
The \(L\) equal site contributions multiplied by the Hamiltonian
prefactor \(g_U/L\) therefore give
\[
 \langle T,H_U\Omega\rangle
  =\frac{g_U}{L}\,L\,\frac{-1}{L^2}
  =-\frac{g_U}{L^2}.
\]
This is nonzero when \(g_U\ne0\), and \(|T\rangle\) is orthogonal
to every zero-species-charge box.  Theorem~\ref{thm:ch21-u-leakage}
then proves the final assertions.  No frozen-margin assumption for
finite density currents has been used.
\end{proof}

\begin{example}[Four-site value]
\label{ex:ch21-u-four-site}
At \(h=2\), \(L=4\), \(\zeta=i\), and \(g_U=1\),
the sea fills \(\{-1,0\}\) in every species.  The witness target
fills \(\{-1,0,1\}\) in each left species and \(\{0\}\) in
each right species.  Equation~\eqref{eq:ch21-u-witness-coefficient}
gives the exact matrix element
\[
 \langle T,H_U\Omega\rangle=-\frac1{16}.
\]
The charge vector \((1,1,-1,-1)\) already suffices to exclude this
ket from the zero-charge box, independently of the excitation budget.
\end{example}

\subsection*{Why a nonzero action need not leak: a Pauli-blocked counterexample}

\begin{example}[Only the inward channel survives]
\label{ex:ch21-u-inward-only}
Take \(h=2\), \(L=4\), and \(g_U\ne0\).  Choose a normalized occupation
ket \(|S\rangle\) with both left species \emph{fully} occupied,
\(\{-1,0,1,2\}\), and both right species occupied only at \(\{-1\}\).
Its relative charges are
\[
  \mathbf N(S)=(2,2,-1,-1).
\]
Use the proper charge box with bounds
\(\mathbf B=(2,2,1,1)\) in the species order
\((L\uparrow,L\downarrow,R\uparrow,R\downarrow)\), and the energy
cutoff \(K=16\).  The unprojected outward sum annihilates \(|S\rangle\):
\[
 O_U(x)|S\rangle=0\qquad\text{for every }x,
\]
because it would attempt to create a fermion in an already full
left species.  On the other hand, the adjoint sum is nonzero and
shifts the charges to \((1,1,0,0)\), which obey the same bounds.

To exhibit a coefficient, let \(|T'\rangle\) be the occupation ket with
left species \(\{0,1,2\}\), right-up \(\{-1,2\}\), and
right-down \(\{-1,0\}\).  The adjoint word removes \(-1\) from
each left species and inserts \(2\) in right-up and \(0\) in
right-down.  Its right-to-left CAR preceding counts are
\(0,3,8,7\), with an even sum \(18\), so the sign is positive.
The Fourier phase is \(\zeta^{(-1-1-2-0)x}=\zeta^{-4x}=1\).
Exactly as in the preceding proof,
\begin{equation}
 \langle T',H_U S\rangle=\frac{g_U}{16}\ne0.
 \label{eq:ch21-u-inward-coefficient}
\end{equation}
For this four-site band, the excitation energy above a fixed-number
sector ground is at most \(n(4-n)\le4\) for a species with \(n\)
occupied modes; summing four species gives \(\mathcal E\le16\).
Thus \emph{every} ket generated from \(|S\rangle\) by the adjoint
Umklapp term belongs to \(\mathcal B_{16,\mathbf B}\).  We have
\(H_U|S\rangle\ne0\) but \((I-P_{16,\mathbf B})H_U|S\rangle=0\).
This proves only that \emph{this state} does not leak; it does not
establish invariance of the whole box under \(H_U\).
\end{example}

\subsection*{What finite leakage does, and does not, establish}

The exact conclusions of this section can now be separated.  The CAR
calculation proves a four-species selection rule and identifies precisely
which charge sectors an Umklapp interaction couples.  Same-sector
first-order compression vanishes, but a charge-bounded or energy-bounded
subspace can fail to be invariant, as witnessed by a nonzero
\emph{outside} matrix element.  The finite compression
\(P_{K,\mathbf B}H_UP_{K,\mathbf B}\) remains a legitimate Hermitian
matrix on the retained carrier; it is simply not the restriction of an
invariant ambient dynamics when the leakage is nonzero.

The finite energy-shell machinery developed in Section~\ref{sec:ch21-finite-effective}
can subsequently be used to organize perturbative block reduction. Its
quadratic Schrieffer--Wolff/Baker--Campbell--Hausdorff identities concern
stated projectors and polynomial coefficients through a fixed order; they do
not automatically provide an analytic conjugation, error estimate, or
RG semigroup.  In particular, removing a shell of \emph{total excitation
energy} is not identical to removing a particular Fourier momentum
mode.  The shell-compression formalism must be instantiated with its
projector and block-diagonality hypotheses before being applied to
\(H_U\).

Finally, a Mott gap would require a spectral statement for a specified
many-body Hamiltonian, and a physically meaningful thermodynamic claim
would require a controlled sequence of systems as \(L\to\infty\),
with a specified density, coupling scaling, and limiting observable.
Likewise the Kosterlitz--Thouless flow equations belong to a
separately constructed continuum effective theory; the presence of a
nonzero finite matrix element \(-g_U/L^2\) neither derives their beta
functions nor proves gap formation.  Any proposed use of bosonized
charge vertices also requires care with the distinction between a
finite-CAR operator, an abstract oscillator representation, and a
coefficientwise formal series.  Chapter~21 remains unimplemented in
Lean at the time of this draft: prior Core support for CAR, charge
sectors and finite projections is not a verification of the new
Umklapp theorems.

\section{Finite Energy-Shell Projection and Effective Hamiltonians}
\label{sec:ch21-finite-effective}

The preceding sections built finite operators before asking which states they
preserve.  In particular, the Umklapp word has a definite change of four
species charges, but this selection rule is not a spectral statement, and its
nonzero action need not preserve a chosen finite-energy carrier.  We now face
a different, constructive question: \emph{if we intentionally retain only a
low-energy part of a finite Hilbert space, what operator describes its
interaction with the discarded states?}  Simply deleting the discarded
states removes the very virtual excursions that can modify low-energy
matrix elements.  A finite Schrieffer--Wolff transformation supplies a
controlled algebraic answer through second order in a coupling parameter.

This section develops the method rather than presupposing familiarity with
renormalization.  The starting objects are ordinary finite-dimensional
Hilbert spaces, orthogonal projections and matrix commutators.  We shall
separate (i) exact partition identities, (ii) a solvable non-resonant matrix
equation, (iii) an \emph{exact identity between polynomials of degree two},
and (iv) approximation of an analytic unitary rotation.  Their hypotheses
are different.  In particular, none of these steps by itself constructs a
charge gap or a continuum renormalization-group (RG) flow.

\subsection{What does an energy shell discard?}
\label{subsec:ch21-energy-shell}

Fix a finite lattice of length $L=2h$ with $h\in\mathbb N_{>0}$, a
positive-mode cutoff $M\geq1$, a set of species labels (in this chapter,
$\{R,L\}\times\{\uparrow,\downarrow\}$), and an admissible integer
species-charge vector $\boldsymbol N$.  The finite CAR Fock carrier will be
denoted $\mathcal F_L$.  We use the occupation-space excitation grading
$\mathsf E\colon\{\text{occupation configurations}\}\to\mathbb N_0$,
relative to the specified charge-sector ground configuration.  The
charge-sector \emph{budget} is the span
\begin{equation}
\mathcal B_{\boldsymbol N,K}
 :=\operatorname{span}\bigl\{\lvert A\rangle:
 \boldsymbol N(A)=\boldsymbol N,\ \mathsf E(A)\leq K\bigr\}
 \subseteq\mathcal F_L,
\label{eq:ch21-budget-definition}
\end{equation}
where $K\in\mathbb N_0$.  We emphasize that \eqref{eq:ch21-budget-definition}
is a span of occupation kets selected by a scalar \emph{grade}; it is not
a vector inequality of the form $H\psi\leq K\psi$.  The allowed grades,
charges and finite-volume margins are those of the earlier energy-budget
construction (see Chapter~\ref{ch:vacuum_budget}).

For $K\geq1$, let $\Pi_{\boldsymbol N,K}$ be the orthogonal coordinate
projection onto this budget and define the two \emph{ambient} projections
\begin{align}
P&:=\Pi_{\boldsymbol N,K-1},\\
Q&:=\Pi_{\boldsymbol N,K}-\Pi_{\boldsymbol N,K-1}
  =\operatorname{proj}_{\mathcal H_{\boldsymbol N,K}},\qquad
\mathcal H_{\boldsymbol N,K}:=
 \operatorname{span}\{\lvert A\rangle:\boldsymbol N(A)=\boldsymbol N,
                                  \mathsf E(A)=K\}.
\label{eq:ch21-shell-projections}
\end{align}
Both projections act on $\mathcal F_L$, even though their ranges are
subspaces of $\mathcal B_{\boldsymbol N,K}$.  Their defining coordinate
supports are disjoint.

\begin{proposition}[Exact finite shell partition]
\label{prop:ch21-shell-partition}
For $K\geq1$ the ambient projections satisfy
\begin{equation}
 P^2=P=P^\dagger,\quad Q^2=Q=Q^\dagger,
 \quad PQ=QP=0,\quad P+Q=\Pi_{\boldsymbol N,K}.
\label{eq:ch21-ambient-partition}
\end{equation}
Consequently, with $\mathcal H:=\mathcal B_{\boldsymbol N,K}$ as the
\emph{new carrier}, the restricted operators obey $P+Q=I_{\mathcal H}$
and give the orthogonal decomposition
$\mathcal H=\operatorname{ran}P\oplus\operatorname{ran}Q$.
\end{proposition}
\begin{proof}
In the orthonormal occupation basis, every coordinate projector acts as
identity on a selected ket and as zero on the others.  Its square therefore
equals itself, and its matrix is real diagonal, hence self-adjoint.  The
supports $\mathsf E\leq K-1$ and $\mathsf E=K$ are disjoint: applying
$PQ$ or $QP$ to each basis ket gives zero.  Their union is exactly the
support $\mathsf E\leq K$, within the fixed charge sector.  Thus their
sum equals $\Pi_{\boldsymbol N,K}$ on every occupation ket, proving
\eqref{eq:ch21-ambient-partition}.  This last projection is the identity
\emph{only after} restricting the carrier to its range.  Orthogonality and
direct-sum decomposition follow by writing every $\psi\in\mathcal H$ as
$\psi=P\psi+Q\psi$ and observing that the ranges intersect trivially.
\end{proof}

\begin{example}[Removing energy two is not removing mode two]
\label{ex:ch21-shell-versus-mode}
Consider the polynomial oscillator carrier with two positive modes,
$X_1$ and $X_2$, weighted by energies $1$ and $2$, respectively.  For
$M\geq2$, the illustrative filtered monomial basis at budget two is
\begin{equation}
 \mathcal B_2^{\mathrm{osc}}
  =\operatorname{span}\{1,X_1,X_1^2,X_2\},\qquad
 \mathcal B_1^{\mathrm{osc}}=\operatorname{span}\{1,X_1\}.
\label{eq:ch21-osc-shell-example}
\end{equation}
The shell removed by $P$ is
$\operatorname{span}\{X_1^2,X_2\}$.  Deleting the \emph{mode} $X_2$
instead would retain $X_1^2$ and yield
$\operatorname{span}\{1,X_1,X_1^2\}$.  These are distinct projections.
Their difference persists in a concrete finite-Fock realization: under
the charge-admissibility and frozen-margin hypotheses required for the
repeated positive mode-one density excitation, the vector
$\rho_{1,\nu}^{\,2}\lvert\boldsymbol N\rangle_0$ is nonzero, has
excitation grade two, is killed by the budget-one projector, and is
fixed by the exact-grade-two shell projector.  Here $\rho_{1,\nu}$
uses the convention that positive $m$ raises the excitation grade.
The needed nonvanishing is a finite CAR/current statement, not an
assumption that the polynomial oscillator is literally the full Fock
carrier.  The supporting Core A10 result gives this witness with
$h\geq2$ and its stated admissibility/margin conditions.
\end{example}

The letter $K$ will continue to denote the \emph{energy budget}, not a
Luttinger parameter.  The shell splitting does not assert that its
grade is the interacting Hamiltonian's spectrum: the grade is a basis
filtration fixed in advance.  A Hamiltonian may mix these spaces, which
is precisely why an effective Hamiltonian is needed.

\subsection{Block matrices and the denominator hypothesis}
\label{subsec:ch21-block-setup}

Work now on an arbitrary finite-dimensional complex Hilbert space
$\mathcal H=\mathcal P\oplus\mathcal Q$, with orthogonal projections
$P$ and $Q=I_{\mathcal H}-P$.  This could be the restricted budget
carrier from Proposition~\ref{prop:ch21-shell-partition}; alternatively it
could be a different explicitly selected finite spectral window.  Let
\begin{equation}
H(t)=H_0+tV,\qquad H_0^\dagger=H_0,\quad V^\dagger=V,
\quad t\in\mathbb R,\quad [P,H_0]=0.
\label{eq:ch21-H-perturbation}
\end{equation}
The last hypothesis means that $H_0$ preserves $\mathcal P$ and
$\mathcal Q$.  It is not automatic for a chosen interacting baseline.
In particular, an interbranch pairing term can mix bare Haldane
partitions: one must diagonalize the proposed $H_0$ on an appropriate
carrier or use a baseline, such as the free Sugawara Hamiltonian, for
which the requisite invariance actually holds.

Insert $I=P+Q$ to write each operator in a $2\times2$ array of blocks:
\begin{equation}
H_0=\begin{pmatrix}A&0\\0&D\end{pmatrix},\qquad
V=\begin{pmatrix}V_{PP}&B\\B^\dagger&V_{QQ}\end{pmatrix},
\quad A=PH_0P\big|_{\mathcal P},\quad D=QH_0Q\big|_{\mathcal Q},
\label{eq:ch21-block-matrices}
\end{equation}
where $B=PVQ\colon\mathcal Q\to\mathcal P$.  Define
\begin{equation}
 V_{\mathrm d}:=PVP+QVQ,\qquad
 V_{\mathrm o}:=PVQ+QVP,\qquad V=V_{\mathrm d}+V_{\mathrm o}.
\label{eq:ch21-diag-offdiag}
\end{equation}
The subscripts refer to this \emph{chosen decomposition}, not to
momentum conservation or locality.  A block-off-diagonal perturbation
moves a state between retained and discarded sectors.  Its second-order
virtual excursion may start and end in $\mathcal P$ even when the first
order transition does not stay there.

Because $H_0$ is self-adjoint and preserves the two summands, we may
choose an orthonormal eigenbasis
$\{\lvert\lambda\rangle\}_{\lambda\in I_P}$ of $A$ and
$\{\lvert\mu\rangle\}_{\mu\in I_Q}$ of $D$,
with \emph{real} eigenvalues $E_\lambda$ and $E_\mu$.  No relation between
the two eigenvalue sets has yet been assumed.  A virtual transition
requires an energy denominator; it is therefore necessary to state the
non-resonance condition rather than dividing by a possibly zero number.

\begin{definition}[Coupling-wise non-resonance]
\label{def:ch21-coupling-nonresonance}
The pair $(H_0,V)$ is \emph{non-resonant across $P/Q$ on its nonzero
couplings} if
\begin{equation}
 \langle\lambda\lvert V\rvert\mu\rangle\ne0
 \quad\Longrightarrow\quad E_\lambda\ne E_\mu,
 \qquad \lambda\in I_P,\ \mu\in I_Q.
\label{eq:ch21-nonresonance}
\end{equation}
Degenerate pairs with zero cross-block matrix element are allowed.
\end{definition}

The stronger assumption
$\operatorname{spec}(A)\cap\operatorname{spec}(D)=\varnothing$
guarantees invertibility of the Sylvester operator introduced below.
The weaker, coupling-wise assumption is sufficient for \emph{existence}
of the particular solution used in this section, but not generally for
uniqueness.  Since there are finitely many coupled pairs, its nonzero
energy differences have a strictly positive minimum when any such
pairs exist; that finite number need not remain bounded away from zero
as $L$ or the budget varies.

\subsection{Solving the Sylvester equation}
\label{subsec:ch21-sylvester}

We seek a first-order change of basis generated by an anti-Hermitian,
block-off-diagonal operator $S_1$:
\begin{equation}
 S_1^\dagger=-S_1,\qquad PS_1P=QS_1Q=0,
 \qquad [S_1,H_0]=-V_{\mathrm o}.
\label{eq:ch21-generator-equation}
\end{equation}
Anti-Hermiticity ensures $e^{tS_1}$ is unitary for real $t$, while the
commutator condition cancels the first-order coupling between
$\mathcal P$ and $\mathcal Q$.  It is important that the commutator is
ordered $[S_1,H_0]=S_1H_0-H_0S_1$; reversing it reverses the
energy-denominator sign.

Write
$S_1=\bigl(\begin{smallmatrix}0&X\\-X^\dagger&0\end{smallmatrix}\bigr)$,
where $X:\mathcal Q\to\mathcal P$.  Multiplying the block matrices in
\eqref{eq:ch21-block-matrices} gives
\begin{equation}
 [S_1,H_0]_{PQ}=XD-AX=-B
 \quad\Longleftrightarrow\quad AX-XD=B.
\label{eq:ch21-sylvester}
\end{equation}
This is a \emph{Sylvester equation}: a linear equation for an unknown
operator $X$ in which fixed operators multiply it on the left and right.
It is not a division by the matrix $A-D$, since the matrices can act on
different spaces and need not commute.

\begin{theorem}[Finite non-resonant Schrieffer--Wolff generator]
\label{thm:ch21-sw-generator}
Under \eqref{eq:ch21-H-perturbation} and
Definition~\ref{def:ch21-coupling-nonresonance}, define the
cross-block matrix entries
\begin{equation}
\langle\lambda\lvert X\rvert\mu\rangle:=
\begin{cases}
\displaystyle\frac{\langle\lambda\lvert V\rvert\mu\rangle}
 {E_\lambda-E_\mu}, &
       \langle\lambda\lvert V\rvert\mu\rangle\ne0,\\[6pt]
0,&\langle\lambda\lvert V\rvert\mu\rangle=0,
\end{cases}
\qquad \lambda\in I_P,\ \mu\in I_Q.
\label{eq:ch21-X-entries}
\end{equation}
Then $S_1=\bigl(\begin{smallmatrix}0&X\\-X^\dagger&0\end{smallmatrix}\bigr)$
is anti-Hermitian, block-off-diagonal, and solves
$[S_1,H_0]=-V_{\mathrm o}$.  The zero choice on uncoupled pairs remains
well defined even when their energies coincide.
\end{theorem}
\begin{proof}
The denominator is nonzero in the first case by
\eqref{eq:ch21-nonresonance}; the second case involves no division.
In the chosen eigenbases,
\begin{align*}
\langle\lambda\lvert AX-XD\rvert\mu\rangle
 &=(E_\lambda-E_\mu)
     \langle\lambda\lvert X\rvert\mu\rangle\\
 &=\langle\lambda\lvert V\rvert\mu\rangle,
\end{align*}
including when the matrix element of $V$ is zero.  Thus $AX-XD=B$
entrywise.  The opposite off-diagonal block of
$[S_1,H_0]$ is the adjoint of the first block, because
$S_1^\dagger=-S_1$ and $H_0^\dagger=H_0$ make $[S_1,H_0]$ Hermitian.
More explicitly, the opposite entry is
\begin{equation*}
 \langle\mu\lvert S_1\rvert\lambda\rangle
 =-\overline{\langle\lambda\lvert S_1\rvert\mu\rangle}
 =\frac{\langle\mu\lvert V\rvert\lambda\rangle}{E_\mu-E_\lambda}
\end{equation*}
when coupled, and zero otherwise.  Both diagonal blocks of
$[S_1,H_0]$ vanish because $H_0$ is block diagonal and $S_1$ is
off diagonal.  Collecting the four blocks proves the commutator
equation; the displayed block form proves anti-Hermiticity directly.
\end{proof}

\begin{remark}[Why the solution need not be unique]
If $E_\lambda=E_\mu$ for an uncoupled cross-block pair, the
corresponding entry of the homogeneous equation
$AX-XD=0$ is unconstrained.  Setting it to zero is a convenient,
explicit choice, not a uniqueness theorem.  If the spectra of $A$ and
$D$ are disjoint, the homogeneous equation has only the zero solution,
and the constructed $X$ is unique.
\end{remark}

\subsection{The ordered second-order transformation}
\label{subsec:ch21-sw-bch}

For any finite matrices $S$ and $T$, write
$\operatorname{ad}_S(T):=[S,T]$ and
$\operatorname{ad}_S^j(T)$ for its $j$-fold iteration.  We rotate the
Hamiltonian according to
\begin{equation}
 \widetilde H(t):=e^{tS_1}(H_0+tV)e^{-tS_1}.
\label{eq:ch21-sw-rotated}
\end{equation}
The orientation is fixed: we conjugate by $e^{+tS_1}$ on the left.
The finite-matrix exponential exists for every real $t$ and the
rotated Hamiltonian is self-adjoint and isospectral to $H(t)$, even
though the compression $P\widetilde H(t)P$ need not reproduce the exact
low-energy spectrum at a finite, non-small $t$.

To obtain the perturbative expansion, differentiate
$e^{tS_1}Te^{-tS_1}$ at $t=0$.  Its successive derivatives are
$\operatorname{ad}_{S_1}^j(T)$.  Hence the ordered polynomial through
degree two is
\begin{align}
B_2(t):={}&H_0+t\bigl(V+[S_1,H_0]\bigr)\notag\\
 &+t^2\left([S_1,V]
        +\frac12[S_1,[S_1,H_0]]\right).
\label{eq:ch21-bch-two-polynomial}
\end{align}
The equality between $B_2$ and the Taylor polynomial is a finite
algebraic fact, irrespective of whether $t$ is small.  The claim that
it approximates $\widetilde H(t)$ needs a separate remainder statement.

\begin{theorem}[Exact compressed degree-two identity]
\label{thm:ch21-sw-compressed}
Assume that $P,Q$ are complementary orthogonal projections on
$\mathcal H$, $H_0$ commutes with $P$, and $S_1$ obeys
\eqref{eq:ch21-generator-equation}.  Then, as an \emph{exact polynomial
identity} on $\mathcal P=\operatorname{ran}P$,
\begin{equation}
 P B_2(t)P
 =PH_0P+tPVP+\frac{t^2}{2}P[S_1,V]P.
\label{eq:ch21-sw-polynomial}
\end{equation}
The right-hand side is the second-order Schrieffer--Wolff effective
Hamiltonian, denoted $H_{\mathrm{eff}}^{[2]}(t)$.  No exponential
appears in the asserted equality.
\end{theorem}
\begin{proof}
Substitute $[S_1,H_0]=-V_{\mathrm o}$ into
\eqref{eq:ch21-bch-two-polynomial}.  The first-order coefficient is
$V-V_{\mathrm o}=V_{\mathrm d}$; in particular,
$PV_{\mathrm d}P=PVP$.  For the second-order term,
\begin{align*}
[S_1,V]+\tfrac12[S_1,[S_1,H_0]]
 &= [S_1,V]-\tfrac12[S_1,V_{\mathrm o}]\\
 &= [S_1,V_{\mathrm d}]+\tfrac12[S_1,V_{\mathrm o}].
\end{align*}
A product of an off-diagonal and a diagonal block operator is
off-diagonal, so $P[S_1,V_{\mathrm d}]P=0$.  Therefore the compressed
coefficient becomes
$\tfrac12P[S_1,V_{\mathrm o}]P$.  Moreover,
$P[S_1,V]P=P[S_1,V_{\mathrm o}]P$ by the same block argument.
The constant term is $PH_0P$, yielding exactly
\eqref{eq:ch21-sw-polynomial}.  Every equality used here is an
ordinary finite-matrix identity, with no perturbative limiting
argument.
\end{proof}

\begin{proposition}[Intermediate-state formula]
\label{prop:ch21-sw-intermediate}
For $\lambda,\lambda'\in I_P$ the second-order correction has matrix
entries
\begin{equation}
\begin{aligned}
\langle\lambda\lvert W^{(2)}\rvert\lambda'\rangle
  &=\frac12\sum_{\substack{\mu\in I_Q\\
       V_{\lambda\mu}V_{\mu\lambda'}\ne0}}
  V_{\lambda\mu}V_{\mu\lambda'}
  \left(\frac{1}{E_\lambda-E_\mu}
        +\frac{1}{E_{\lambda'}-E_\mu}\right),\\
W^{(2)}&:=\frac12P[S_1,V]P,
\qquad V_{ab}:=\langle a\lvert V\rvert b\rangle.
\end{aligned}
\label{eq:ch21-sw-intermediate-sum}
\end{equation}
Only nonzero coupled terms are included in the displayed sum, so its
denominators are defined under \eqref{eq:ch21-nonresonance}.  In
particular,
\begin{equation}
\langle\lambda\lvert W^{(2)}\rvert\lambda\rangle
 =\sum_{\substack{\mu\in I_Q\\V_{\lambda\mu}\ne0}}
 \frac{|V_{\lambda\mu}|^2}{E_\lambda-E_\mu}.
\label{eq:ch21-sw-diagonal-shift}
\end{equation}
\end{proposition}
\begin{proof}
Since $S_1$ is off diagonal, the $PP$ block of $S_1V$ is
$S_{1,PQ}V_{QP}$ and that of $VS_1$ is $V_{PQ}S_{1,QP}$.
Insert the orthonormal resolution $I_{\mathcal Q}=\sum_{\mu\in I_Q}
\lvert\mu\rangle\langle\mu\rvert$ to get
\begin{equation*}
\langle\lambda\lvert[S_1,V]\rvert\lambda'\rangle
=\sum_{\mu\in I_Q}\bigl((S_1)_{\lambda\mu}
 V_{\mu\lambda'}-V_{\lambda\mu}(S_1)_{\mu\lambda'}\bigr).
\end{equation*}
Theorem~\ref{thm:ch21-sw-generator} gives
$(S_1)_{\lambda\mu}=V_{\lambda\mu}/(E_\lambda-E_\mu)$
on coupled pairs and
$(S_1)_{\mu\lambda'}=-V_{\mu\lambda'}/(E_{\lambda'}-E_\mu)$.
Terms with vanishing couplings contribute zero by definition.  This
yields \eqref{eq:ch21-sw-intermediate-sum}.  Taking
$\lambda'=\lambda$ and using Hermiticity of $V$ proves
\eqref{eq:ch21-sw-diagonal-shift}.
\end{proof}

The denominator in \eqref{eq:ch21-sw-diagonal-shift} explains the
physical interpretation: the retained state makes a virtual excursion
to a discarded eigenstate and returns.  In a strictly ordered baseline
spectrum $E_\lambda<E_\mu$ for every coupled pair, each diagonal
second-order shift is nonpositive.  Without that ordering, cancellations
are possible.  Thus even the \emph{sign} of the correction is not
determined by the mere existence of a $P/Q$ splitting.

\subsection{What is the remainder, exactly?}
\label{subsec:ch21-sw-remainder}

There are two different uses of the notation $O(t^3)$.  In a formal
power series ring $\operatorname{End}(\mathcal H)[[t]]$, the statement
\begin{equation}
P\widetilde H(t)P\equiv H_{\mathrm{eff}}^{[2]}(t)
 \pmod{t^3}
\label{eq:ch21-sw-formal-mod}
\end{equation}
means \emph{coefficientwise} equality in degrees zero, one and two.
It does not evaluate an operator at a numerical $t$, and it does not
by itself bound the discarded terms.  For the finite analytic
exponential we can additionally obtain an explicit bound; this is
an elementary finite-matrix consequence, stated here separately from
the supplied A10 Lean theorem.

\begin{proposition}[Finite analytic Taylor remainder]
\label{prop:ch21-sw-analytic-remainder}
Under the hypotheses of Theorem~\ref{thm:ch21-sw-compressed}, for real
$t$ let
\begin{equation}
 R_3(t):=\widetilde H(t)-B_2(t).
\label{eq:ch21-rem-definition}
\end{equation}
Then the exact compressed analytic Hamiltonian obeys
\begin{equation}
 P\widetilde H(t)P=H_{\mathrm{eff}}^{[2]}(t)+PR_3(t)P,
\label{eq:ch21-sw-exact-plus-remainder}
\end{equation}
where, writing $\mathcal U_r(T)=e^{rS_1}Te^{-rS_1}$,
\begin{align}
R_3(t)={}&\frac{t^3}{2}\int_0^1(1-u)^2
  \mathcal U_{ut}\bigl(\operatorname{ad}_{S_1}^3(H_0)\bigr)\,du\notag\\
 &+t^3\int_0^1(1-u)
  \mathcal U_{ut}\bigl(\operatorname{ad}_{S_1}^2(V)\bigr)\,du.
\label{eq:ch21-sw-integral-remainder}
\end{align}
For the operator norm, using $S_1^\dagger=-S_1$,
\begin{equation}
\|PR_3(t)P\|\leq\|R_3(t)\|
\leq |t|^3\left(
\frac{\|\operatorname{ad}_{S_1}^3(H_0)\|}{6}
+\frac{\|\operatorname{ad}_{S_1}^2(V)\|}{2}\right).
\label{eq:ch21-sw-norm-bound}
\end{equation}
The bound is finite-volume and its constant depends on the chosen
carrier, $H_0,V$ and the selected generator $S_1$.
\end{proposition}
\begin{proof}
For a fixed $T$, put $F_T(t)=e^{tS_1}Te^{-tS_1}$.
Differentiating the matrix exponential gives
$F_T^{(j)}(r)=\mathcal U_r(\operatorname{ad}_{S_1}^jT)$.
Taylor's theorem with integral remainder through order two for
$F_{H_0}$ gives
\begin{equation*}
F_{H_0}(t)=H_0+t\operatorname{ad}_{S_1}(H_0)
+\frac{t^2}{2}\operatorname{ad}_{S_1}^2(H_0)
+\frac{t^3}{2}\int_0^1(1-u)^2
  \mathcal U_{ut}(\operatorname{ad}_{S_1}^3H_0)\,du.
\end{equation*}
Taylor's theorem through order one for $F_V$ gives
\begin{equation*}
F_V(t)=V+t\operatorname{ad}_{S_1}(V)
+t^2\int_0^1(1-u)
 \mathcal U_{ut}(\operatorname{ad}_{S_1}^2V)\,du.
\end{equation*}
Add $F_{H_0}(t)+tF_V(t)$ and compare with
\eqref{eq:ch21-bch-two-polynomial}; this proves
\eqref{eq:ch21-sw-integral-remainder} and then
\eqref{eq:ch21-sw-exact-plus-remainder}.
For real $t$ and $u$, $e^{utS_1}$ is unitary, so
$\|\mathcal U_{ut}(T)\|=\|T\|$.  Apply the triangle inequality and
$\int_0^1(1-u)^2du=1/3$ and
$\int_0^1(1-u)du=1/2$ to obtain
\eqref{eq:ch21-sw-norm-bound}.  Finally, orthogonal projections are
contractions, hence $\|PR_3P\|\leq\|R_3\|$.
\end{proof}

\begin{remark}[A finite error is not a uniform scaling estimate]
The right side of \eqref{eq:ch21-sw-norm-bound} may increase with
lattice size or as the non-resonant denominators become small.  It is
a genuine estimate at fixed finite carrier, but not a proof of an
$L\to\infty$ limit or a scale-independent perturbative expansion.
The coefficientwise polynomial result of the supplied Core A10 module
should not be reported as formal verification of this analytic bound.
\end{remark}

\subsection{A two-level calculation, including what fails at finite coupling}
\label{subsec:ch21-two-level}

\begin{example}[Explicit Schrieffer--Wolff rotation]
\label{ex:ch21-sw-two-level}
Let $\mathcal H=\mathbb C^2$, $P=\operatorname{diag}(1,0)$ and
$Q=\operatorname{diag}(0,1)$.  Take
\begin{equation}
 H_0=\begin{pmatrix}0&0\\0&1\end{pmatrix},\qquad
 V=\begin{pmatrix}0&1\\1&0\end{pmatrix}.
\label{eq:ch21-sw-2x2-input}
\end{equation}
The low state has energy zero and the high state energy one.  The
Sylvester denominator is $0-1=-1$, so
\begin{equation}
 S_1=\begin{pmatrix}0&-1\\1&0\end{pmatrix},\qquad
 [S_1,H_0]=-V,
 \quad[S_1,V]=\begin{pmatrix}-2&0\\0&2\end{pmatrix}.
\label{eq:ch21-sw-2x2-generator}
\end{equation}
The effective Hamiltonian is therefore the scalar operator
\begin{equation}
 H_{\mathrm{eff}}^{[2]}(t)=-t^2 P
 \quad\text{on }\operatorname{ran}P.
\label{eq:ch21-sw-2x2-effective}
\end{equation}
We can test exactly which part of this assertion is an approximation.
Since $S_1^2=-I$, the rotation is
\begin{equation}
 e^{tS_1}=\begin{pmatrix}\cos t&-\sin t\\
                                   \sin t&\cos t\end{pmatrix}.
\end{equation}
Direct multiplication gives the \emph{exact} compressed entry
\begin{align}
 (P\widetilde H(t)P)_{11}
 &=\sin^2t-t\sin(2t)
   =-t^2+t^4+O(t^6),
\label{eq:ch21-sw-2x2-rotated}\end{align}
and the off-diagonal entry
\begin{equation}
 (\widetilde H(t))_{12}
 =t\cos(2t)-\tfrac12\sin(2t)
 =-\tfrac43t^3+O(t^5).
\label{eq:ch21-sw-2x2-off}
\end{equation}
Thus the first-order mixing is canceled but the chosen rotation is
not an exact block diagonalization for arbitrary $t$.  Independently,
the exact eigenvalues of $H_0+tV$ are
\begin{equation}
 E_\pm(t)=\frac{1\pm\sqrt{1+4t^2}}2.
\label{eq:ch21-sw-2x2-eigvals}
\end{equation}
At $t=2/3$, the lower eigenvalue is $-1/3$ (with eigenvector
$(-2,1)$), whereas \eqref{eq:ch21-sw-2x2-effective} gives $-4/9$.
The two values differ because a second-order effective operator is
not an exact finite-coupling spectral identification.
\end{example}

The same caution applies when the quadratic correction happens to
vanish.  A nonzero perturbation can couple retained and discarded
states through \emph{different} intermediate energies with cancelling
denominators.

\begin{example}[Zero second-order correction with nonzero leakage]
\label{ex:ch21-sw-three-level}
Take $\mathcal P=\operatorname{span}\{e_1\}$,
$\mathcal Q=\operatorname{span}\{e_2,e_3\}$ and
\begin{equation}
 H_0=\operatorname{diag}(0,-1,1),\quad
 V=\begin{pmatrix}0&1&1\\1&0&0\\1&0&0\end{pmatrix},\quad
 S_1=\begin{pmatrix}0&1&-1\\-1&0&0\\1&0&0\end{pmatrix}.
\label{eq:ch21-sw-three-input}
\end{equation}
Here $PVQ\ne0$ and $[S_1,H_0]=-V$, but
\begin{equation}
\langle e_1\lvert W^{(2)}\rvert e_1\rangle
 =\frac{|V_{12}|^2}{0-(-1)}
  +\frac{|V_{13}|^2}{0-1}=1-1=0.
\label{eq:ch21-sw-three-cancel}
\end{equation}
A vanishing compressed second-order commutator therefore does
\emph{not} imply $PVQ=QVP=0$.  Conversely, if both off-block
couplings vanish, choosing $S_1=0$ gives the exact identity-rotation
compression $P(H_0+tV)P=PH_0P+tPVP$ and zero correction.  That
forward implication does not require a non-resonance hypothesis.
\end{example}

\subsection{Scope of the finite construction and the remaining chapter obligations}
\label{subsec:ch21-effective-scope}

The discrete procedure has three mathematically independent inputs:
(i) a precisely selected coordinate budget and its orthogonal shell,
(ii) an invariant baseline $H_0$ and a cross-block perturbation $V$,
and (iii) control of coupled energy denominators.  Earlier sections
provided exact finite CAR words, charge-selection rules, and conditional
leakage witnesses for spinful perturbations.  Those statements do not
by themselves furnish the spectral decomposition, denominator bounds,
or block invariance needed to insert a particular interacting
Hamiltonian into the Schrieffer--Wolff theorem.

The separately verified finite support results cover the non-resonant
commutator equation and skew-Hermiticity of the constructed generator,
two- and three-level matrix witnesses, the compressed \emph{degree-two
BCH polynomial} identity, the decomposition of shell projections and
their orthogonality and adjoint properties, and a nonzero repeated-mode-one
shell witness. One projector identity is $P+Q=\Pi_{\boldsymbol N,K}$ on the
ambient Fock carrier, whereas the generic compressed BCH theorem assumes
$P+Q=I$ on its whole carrier. Restricting to
$\mathcal B_{\boldsymbol N,K}$ and transferring an actual model's
operators and hypotheses is a separate bridge, not established merely
by placing the two statements next to each other.  These finite support
proofs do not constitute formal verification of the present Chapter~21
manuscript.

The remaining constructive tasks include proving, for any chosen
spinful Hamiltonian, that its finite blocks act on the proposed
carriers; checking charge-sector and energy-margin constraints at
each intermediate application; constructing a valid $H_0$ and its
adapted eigenbasis or another certified Sylvester solver; and proving
any claimed uniform denominator or remainder bounds.  The
coupling-wise condition is conditional, and a vanishing second-order
term is not a substitute for checking it.  Furthermore, a Mott or
spin gap would require an independently defined spectral gap and a
proof of its behavior under relevant size or coupling limits.
A continuum RG flow requires a scale-changing construction and a
comparison between effective theories at successive scales; energy
shell projection at one finite $K$ is not such a construction.
Finally, an isolated finite Hermitian Hamiltonian has a real discrete
spectrum, so scattering lifetimes or dissipative flow cannot be
inferred solely from the matrix identities above.

What has been established here is therefore deliberately narrower
and more precise: a well-defined \emph{finite energy-shell filter}, a
non-resonant \emph{first-order unitary generator}, an exact
\emph{second-order compressed polynomial identity}, and a finite
analytic remainder estimate with explicit hypotheses.  These results
provide the operator-level starting point for a future controlled
effective-Hamiltonian analysis, without prematurely identifying its
outcome with a charge gap or a continuum renormalization-group law.

